% ===================================================================
%  AntennaForge — User Guide and Architecture Design Document
%  Auto-generated comprehensive reference
% ===================================================================
\documentclass[11pt,a4paper,twoside]{article}

% ── Packages ──────────────────────────────────────────────────────
\usepackage[utf8]{inputenc}
\usepackage[T1]{fontenc}
\usepackage{lmodern}
\usepackage[margin=1in]{geometry}
\usepackage{graphicx}
\usepackage{xcolor}
\usepackage{hyperref}
\usepackage{booktabs}
\usepackage{longtable}
\usepackage{tabularx}
\usepackage{array}
\usepackage{enumitem}
\usepackage{amsmath,amssymb}
\usepackage{listings}
\usepackage{fancyhdr}
\usepackage{titlesec}
\usepackage{tocloft}
\usepackage{float}
\usepackage{caption}
\usepackage{multirow}
\usepackage{parskip}

% ── Colour Definitions ───────────────────────────────────────────
\definecolor{afblue}{HTML}{3B82F6}
\definecolor{afdark}{HTML}{0F172A}
\definecolor{afgray}{HTML}{475569}
\definecolor{aflight}{HTML}{F8FAFC}
\definecolor{afgreen}{HTML}{16A34A}
\definecolor{afred}{HTML}{DC2626}
\definecolor{aforange}{HTML}{D97706}
\definecolor{afcode}{HTML}{1E293B}

% ── Hyperref Setup ───────────────────────────────────────────────
\hypersetup{
    colorlinks=true,
    linkcolor=afblue,
    urlcolor=afblue,
    citecolor=afblue,
    pdftitle={AntennaForge User Guide and Architecture},
    pdfauthor={AntennaForge Project},
    pdfsubject={Antenna Pattern Generation, Analysis, and Reporting},
}

% ── Code Listing Style ───────────────────────────────────────────
\lstdefinestyle{afstyle}{
    backgroundcolor=\color{aflight},
    basicstyle=\ttfamily\small,
    keywordstyle=\color{afblue}\bfseries,
    commentstyle=\color{afgray}\itshape,
    stringstyle=\color{afgreen},
    frame=single,
    rulecolor=\color{afgray!30},
    breaklines=true,
    breakatwhitespace=true,
    tabsize=4,
    showstringspaces=false,
    numbers=left,
    numberstyle=\tiny\color{afgray},
}
\lstset{style=afstyle}

% ── Header / Footer ─────────────────────────────────────────────
\pagestyle{fancy}
\fancyhf{}
\fancyhead[LE,RO]{\thepage}
\fancyhead[LO]{\nouppercase{\leftmark}}
\fancyhead[RE]{AntennaForge User Guide}
\renewcommand{\headrulewidth}{0.4pt}

% Fix headheight warning
\setlength{\headheight}{13.59999pt}

% ── Section Formatting ───────────────────────────────────────────
\titleformat{\section}{\Large\bfseries\color{afdark}}{}{0em}{\thesection\quad}
\titleformat{\subsection}{\large\bfseries\color{afblue}}{}{0em}{\thesubsection\quad}
\titleformat{\subsubsection}{\normalsize\bfseries\color{afgray}}{}{0em}{\thesubsubsection\quad}

% ── Custom column types ──────────────────────────────────────────
\newcolumntype{L}[1]{>{\raggedright\arraybackslash}p{#1}}
\newcolumntype{C}[1]{>{\centering\arraybackslash}p{#1}}

% ── Document Begin ───────────────────────────────────────────────
\begin{document}

% ===================================================================
%  TITLE PAGE
% ===================================================================
\begin{titlepage}
\centering
\vspace*{3cm}
{\Huge\bfseries\color{afdark} AntennaForge\\[0.3em]}
{\LARGE\color{afblue} User Guide \& Architecture Design Document\\[2em]}
{\Large Version 1.0\\[1em]}
{\large\color{afgray} Comprehensive Reference for All Features,\\
Fields, Controls, and Internal Architecture\\[3em]}
\vfill
{\large\today}
\end{titlepage}

% ===================================================================
%  TABLE OF CONTENTS
% ===================================================================
\newpage
\tableofcontents
\newpage

% ===================================================================
%  PART I — INTRODUCTION AND OVERVIEW
% ===================================================================
\newpage
\section{Introduction}
\label{sec:introduction}

\textbf{Disclaimer:} AntennaForge generates \emph{synthetic} radiation patterns based on analytical mathematical models (e.g., Gaussian, $\cos^n$, sinc, Jinc). These are idealised approximations designed for system-level simulation. They are \textbf{not} derived from electromagnetic simulation of physical geometry (Method of Moments, FEM, FDTD) and may not capture near-field effects, complex structural scattering, or specific hardware imperfections unless explicitly modelled via the available physics features.

AntennaForge is a Python-based antenna radiation pattern generator, analyser, and reporting tool. It produces realistic two-dimensional gain patterns in CSV format for a wide range of antenna types, supports frequency-dependent behaviour and numerous physics-based features, and provides a complete suite of visualisation, analysis, arithmetic, rotation, interpolation, link budget, and reporting capabilities.

AntennaForge operates in three modes:
\begin{enumerate}[label=\arabic*.]
    \item \textbf{Graphical User Interface (GUI)} — A modern desktop application built with customtkinter, featuring nine tabbed pages with full widget-based control.
    \item \textbf{Interactive CLI} — A terminal-based menu-driven interface with state-machine navigation and GoBack support.
    \item \textbf{Command-Line Interface (CLI)} — Direct invocation via command-line arguments for scriptable, headless operation.
\end{enumerate}

\subsection{Key Capabilities}
\begin{itemize}
    \item Seven built-in antenna models: LPDA, Omni, Panel, Dish, Horn, Monopole, and Phased Array.
    \item Multi-frequency pattern generation with linear or logarithmic spacing.
    \item Nine configurable physics features: frequency-dependent beamwidth (azimuth and elevation independently), Taylor sidelobes, pattern breakup, ground reflection, cross-polarisation, VSWR rolloff, asymmetry/squint, and gain-beamwidth coupling.
    \item Pattern arithmetic (add, subtract, multiply, scale, max envelope, average).
    \item Three-axis pattern rotation (yaw, pitch, roll) with bilinear/bicubic regridding.
    \item Frequency interpolation (two-file, multi-file) and angular resampling.
    \item EIRP calculation, power density computation, and free-space path loss.
    \item Full link budget analysis (boresight margin and angular map modes).
    \item Sidelobe mask compliance testing against ITU-R S.580-6, ITU-R S.465-6, and MIL-STD standards.
    \item PDF and \LaTeX\ report generation with 13 configurable sections, cover pages, and classification markings.
    \item Batch processing from JSON recipe files with dry-run validation.
    \item Animated frequency-sweep GIFs (heatmap and polar styles).
\end{itemize}

\newpage
\subsection{System Requirements}
\begin{itemize}
    \item Python 3.9 or later (3.10+ recommended).
    \item \textbf{Required:} None — AntennaForge can run with the Python standard library alone (pure-Python fallback).
    \item \textbf{Optional (recommended):}
    \begin{itemize}
        \item \texttt{numpy} — Vectorised computation (10--100$\times$ faster pattern generation).
        \item \texttt{scipy} — Bicubic interpolation, Bessel functions for Dish patterns, advanced resampling.
        \item \texttt{matplotlib} — All plotting and report generation.
        \item \texttt{Pillow} — Image handling for the GUI plot viewer.
        \item \texttt{customtkinter} — GUI framework (required only for GUI mode).
        \item \texttt{python-pptx} — PowerPoint report generation.
    \end{itemize}
\end{itemize}

\subsection{Installation}
\begin{lstlisting}[language=bash,caption={Standard installation}]
pip install numpy scipy matplotlib Pillow customtkinter
\end{lstlisting}

An offline installation mode is also supported via bundled wheel files in the \texttt{wheels/} directory. See \texttt{PORTABLE\_SETUP.md} for details.

\subsection{Launching AntennaForge}
\begin{lstlisting}[language=bash,caption={Launch modes}]
# GUI mode
python AntennaForge.py

# Interactive CLI mode
python -m antenna_tool

# Direct CLI generation
python -m antenna_tool --type LPDA --max-gain 10 --az-bw 65 --el-bw 70 \
    --f-min 30 --f-max 300 --slices 10 --output ./patterns
\end{lstlisting}

% ===================================================================
%  PART II — ARCHITECTURE OVERVIEW
% ===================================================================
\newpage
\section{Architecture Overview}
\label{sec:architecture}

\subsection{Directory Structure}
The AntennaForge codebase is organised into the following packages:

\begin{longtable}{L{3.5cm} L{10cm}}
\toprule
\textbf{Package / File} & \textbf{Purpose} \\
\midrule
\endhead
	exttt{reporting/} & PDF report generator (\texttt{report.py}), \LaTeX\ report generator (\texttt{latex\_report.py}), and PowerPoint report generator (\texttt{pptx\_report.py}). The \LaTeX\ generator saves plot images into \texttt{images/} subfolders alongside the \texttt{.tex} file; the PDF generator renders directly into the PDF; the PPTX generator builds slide decks (requires \texttt{python-pptx}). \\
\texttt{AntennaForge.py} & GUI launcher (launches the customtkinter desktop application). \\
\texttt{\_\_main\_\_.py} & Package entry point for \texttt{python -m antenna\_tool} (delegates to \texttt{core/cli.py}). \\
\texttt{core/} & Engine, pattern math, I/O, operations, rotation, interpolation, EIRP, link budget, sidelobe masks, batch processing, parameter fitting, CLI, logging, and dependency checking. \\
\texttt{antennas/} & Base antenna class, antenna registry, and seven antenna model sub-packages (array, dish, horn, lpda, monopole, omni, panel). \\
\texttt{analysis/} & Pattern analysis and statistics (beamwidth measurement, sidelobe detection, symmetry, comparison). \\
\texttt{graphing/} & Matplotlib-based plotting: single-file plots, multi-file plots, animation/GIF generation, and live preview plots for the GUI. \\
\texttt{gui/} & customtkinter desktop application: app shell, theme engine, user preferences, reusable widgets, and nine page modules. \\
\texttt{ui/} & Terminal-mode UI: interactive menus, prompts, and feature configuration wizards. \\
\texttt{scripts/} & Utility scripts: offline wheel download, installation, diff tables, and dataset batch generation. \\
\texttt{tests/} & pytest test suite. \\
\texttt{Docs/} & Documentation (this document). \\
\bottomrule
\end{longtable}

\subsection{Data Flow}
The central data flow of AntennaForge is:
\begin{enumerate}
    \item \textbf{Configuration} — The user specifies antenna type, gain, beamwidths, frequency band, features, and output settings. This is stored as a Python dictionary conforming to the schema in \texttt{config.py}.
    \item \textbf{Generation} — The engine (\texttt{core/engine.py}) iterates over frequency slices. For each frequency, it computes a 2D gain grid over all azimuth/elevation angles using the selected antenna model and all enabled physics features. The result is written to CSV.
    \item \textbf{Post-processing} — The user may apply arithmetic operations, rotation, interpolation, or resampling to the generated CSVs.
    \item \textbf{Analysis} — Pattern CSVs are loaded and analysed for beamwidth, sidelobes, symmetry, and mask compliance.
    \item \textbf{Visualisation} — Patterns are plotted as heatmaps, cuts, polar plots, 3D surfaces, gain-vs-frequency trends, or animated GIFs. Plot images are saved in \texttt{images/<plot\_type>/} subfolders alongside the source CSVs.
    \item \textbf{Reporting} — A complete PDF or \LaTeX\ report is generated from a directory of pattern CSVs. The PDF report renders all figures in-memory via \texttt{PdfPages}. The \LaTeX\ report generates PNG images into \texttt{images/} subfolders (\texttt{heatmaps/}, \texttt{polar/}, \texttt{3d\_surface/}, \texttt{xpol/}) and references them via \texttt{\textbackslash includegraphics}.
\end{enumerate}

\subsection{CSV File Format}
\label{sec:csvformat}
All patterns are stored in a simple CSV format:
\begin{itemize}
    \item \textbf{Row 0 (header):} The first cell is \texttt{0} (placeholder), followed by azimuth angles in degrees (e.g., \texttt{0, 1, 2, \ldots, 359}).
    \item \textbf{Rows 1--$N$:} Each row begins with an elevation angle (e.g., \texttt{-90, -89, \ldots, 90}), followed by gain values in dBi for each azimuth at that elevation.
\end{itemize}

File naming convention: \texttt{<Type>\_<Freq>MHz\_<Pol>.csv}, for example: \texttt{LPDA\_150p00MHz\_copol.csv}. Cross-polarisation files use \texttt{\_xpol.csv} suffix.

\subsection{Plot Image Directory Structure}
\label{sec:plotimages}
Plot images are organised into subfolders by plot type. The \textbf{graphing} modules and the \textbf{\LaTeX\ report generator} each create an \texttt{images/} tree:

\begin{itemize}
    \item \texttt{images/heatmaps/} — Heatmap + az/el cuts combo plots.
    \item \texttt{images/cuts/} — Principal-plane cut plots (graphing module only).
    \item \texttt{images/polar/} — Polar pattern plots.
    \item \texttt{images/3d\_surface/} — 3-D surface plots.
    \item \texttt{images/xpol/} — Cross-polarisation heatmaps (\LaTeX\ report only).
    \item \texttt{images/gain\_vs\_freq/} — Gain-vs-frequency trend (graphing module).
    \item \texttt{images/bw\_vs\_freq/} — Beamwidth-vs-frequency trend (graphing module).
    \item \texttt{images/overlay\_cuts/} — Overlay of all frequency cuts (graphing module).
    \item Global trend plots (\texttt{gain\_vs\_freq.png}, \texttt{bw\_vs\_freq.png}, \texttt{overlay\_cuts.png}) are saved directly to \texttt{images/} by the \LaTeX\ report generator.
\end{itemize}

The PDF report generator (\texttt{reporting/report.py}) renders all figures directly into the PDF via Matplotlib's \texttt{PdfPages} backend and does not write any intermediate PNG files.

\subsection{Configuration System}
\label{sec:config}
The \texttt{DEFAULT\_CONFIG} dictionary in \texttt{config.py} defines every parameter with its default value. Antenna-specific parameters (e.g., \texttt{dish\_efficiency}, \texttt{element\_length\_wavelengths}) are merged at runtime from each antenna model's \texttt{default\_params()} method.

Configuration validation is performed by \texttt{validate\_config()}, which checks:
\begin{itemize}
    \item Required key presence and types.
    \item Numeric ranges (e.g., gain: $-30$ to $60$\,dBi; beamwidth: $1^\circ$ to $360^\circ$; frequency: $0.001$ to $10^6$\,MHz).
    \item Cross-field constraints (e.g., $f_{\min} \leq f_{\max}$).
    \item Valid enumeration values for spacing (\texttt{linear}, \texttt{log}) and polarisation (\texttt{vertical}, \texttt{horizontal}, \texttt{custom}).
    \item Feature sub-dictionary integrity.
    \item Antenna type existence in the registry.
\end{itemize}

% ===================================================================
%  PART III — ANTENNA MODELS
% ===================================================================
\newpage
\section{Antenna Models}
\label{sec:antennas}

AntennaForge includes seven built-in antenna models. Each model inherits from the abstract base class \texttt{AntennaBase} and implements its own gain computation formula. All models are registered in the global \texttt{ANTENNA\_REGISTRY} via the \texttt{register()} mechanism in \texttt{antennas/\_\_init\_\_.py}.

\subsection{Base Antenna Class}
\label{sec:antennabase}
The \texttt{AntennaBase} abstract class (\texttt{antennas/base.py}) defines the interface:
\begin{itemize}
    \item \texttt{name} (abstract property) — display/registry name (e.g., \texttt{"LPDA"}).
    \item \texttt{has\_az\_pattern} (property) — \texttt{True} if the antenna shapes the azimuth pattern (defaults to \texttt{True}; overridden to \texttt{False} for Omni and Monopole).
    \item \texttt{default\_params()} (abstract) — returns a dictionary of antenna-specific default parameters.
    \item \texttt{configure(cfg)} (abstract) — interactive prompts for antenna-specific settings (CLI mode).
    \item \texttt{validate\_config(cfg)} — returns a list of warning strings for unusual configurations.
    \item \texttt{compute\_point\_gain($\ldots$)} (abstract) — returns the co-pol gain (linear scale) at one (azimuth, elevation) point.
    \item \texttt{file\_prefix} (property) — prefix for generated CSV filenames.
\end{itemize}

% ── LPDA ─────────────────────────────────────────────────────────
\subsection{LPDA (Log-Periodic Dipole Array)}
\label{sec:lpda}

\textbf{Registry name:} \texttt{LPDA} \quad\textbar\quad \textbf{Has azimuth pattern:} Yes

\subsubsection{Default Parameters}
\begin{tabular}{lll}
\toprule
\textbf{Parameter} & \textbf{Default} & \textbf{Description} \\
\midrule
\texttt{ftb\_ratio\_db} & 15.0 & Front-to-back ratio in dB \\
\bottomrule
\end{tabular}

\subsubsection{Mathematical Model}
The LPDA uses an elliptical Gaussian main beam model:
\begin{equation}
G_{\text{front}} = G_{\text{peak}} \cdot 0.5^{r^2}
\end{equation}
where the radial distance $r$ is computed as:
\begin{equation}
r^2 = \left(\frac{\mathrm{az}_{\text{eff}}}{\mathrm{az}_{\text{hw}}}\right)^2 + \left(\frac{\mathrm{el}_{\text{eff}}}{\mathrm{el}_{\text{hw}}}\right)^2
\end{equation}
with $\mathrm{az}_{\text{hw}}$ and $\mathrm{el}_{\text{hw}}$ being the half-power half-widths in azimuth and elevation respectively. Sidelobes are added via an elliptical radial envelope. The back lobe is modelled as:
\begin{equation}
G_{\text{back}} = \frac{G_{\text{peak}}}{\mathrm{FTB}_{\text{linear}}} \cdot |\cos(\mathrm{el})|
\end{equation}
The total gain is the maximum of front and back contributions, ensuring a smooth transition with no dead zones at $\pm 90^\circ$.

\subsubsection{Validation Warnings}
\begin{itemize}
    \item Front-to-back ratio $\leq 0$\,dB.
    \item Azimuth beamwidth $\geq 360^\circ$.
    \item Gain $\leq 0$\,dBi.
\end{itemize}

% ── Omni ─────────────────────────────────────────────────────────
\subsection{Omnidirectional (Omni)}
\label{sec:omni}

\textbf{Registry name:} \texttt{Omni} \quad\textbar\quad \textbf{Has azimuth pattern:} No (360° uniform in azimuth)

\subsubsection{Default Parameters}
No antenna-specific parameters (empty dictionary). Azimuth beamwidth is forced to $360^\circ$.

\subsubsection{Mathematical Model}
The omni pattern has uniform azimuth response and shapes only the elevation:
\begin{equation}
G(\theta_{\text{el}}) = G_{\text{peak}} \cdot |\cos(\theta_{\text{el}})|^{n_{\text{el}}}
\end{equation}
where $n_{\text{el}}$ is derived from the elevation beamwidth via the \texttt{bw\_to\_exponent()} function:
\begin{equation}
n = \frac{-3}{10 \cdot \log_{10}\!\bigl(\cos(\mathrm{BW}/2)\bigr)}
\end{equation}
A sidelobe envelope is added in the elevation plane.

\subsubsection{Validation Warnings}
\begin{itemize}
    \item Azimuth beamwidth $\neq 360^\circ$.
\end{itemize}

% ── Panel ────────────────────────────────────────────────────────
\subsection{Panel (Sector Antenna)}
\label{sec:panel}

\textbf{Registry name:} \texttt{Panel} \quad\textbar\quad \textbf{Has azimuth pattern:} Yes

\subsubsection{Default Parameters}
\begin{tabular}{lll}
\toprule
\textbf{Parameter} & \textbf{Default} & \textbf{Description} \\
\midrule
\texttt{ftb\_ratio\_db} & 20.0 & Front-to-back ratio in dB \\
\texttt{mechanical\_tilt\_deg} & 0.0 & Mechanical uptilt/downtilt in degrees (positive = down) \\
\texttt{max\_gain\_dbi} & 16.0 & Peak gain in dBi \\
\texttt{az\_beamwidth\_deg} & 65.0 & Azimuth 3\,dB beamwidth \\
\texttt{el\_beamwidth\_deg} & 10.0 & Elevation 3\,dB beamwidth \\
\texttt{f\_min\_mhz} & 700.0 & Lower band edge in MHz \\
\texttt{f\_max\_mhz} & 2700.0 & Upper band edge in MHz \\
\bottomrule
\end{tabular}

\textbf{Note:} \texttt{mechanical\_tilt\_deg} replaces the legacy \texttt{panel\_downtilt\_deg} parameter. Both are accepted for backward compatibility, but \texttt{mechanical\_tilt\_deg} takes precedence. This parameter is also available for Horn, Dish, and Array types.

\subsubsection{Mathematical Model}
The Panel antenna uses a sinc-squared model in both planes:
\begin{equation}
G(\theta_{\text{az}}, \theta_{\text{el}}) = G_{\text{peak}} \cdot \operatorname{sinc}^2(u_{\text{az}}) \cdot \operatorname{sinc}^2(u_{\text{el}})
\end{equation}
where:
\begin{equation}
u = 1.3916 \cdot \frac{\sin(\theta)}{\sin(\theta_{3\text{dB}}/2)}
\end{equation}
Mechanical tilt is applied by offsetting the elevation angle before gain computation: $\theta'_{el} = \theta_{el} - \text{tilt}$. The back lobe is weighted by the front-to-back ratio and elevation.

\subsubsection{Validation Warnings}
\begin{itemize}
    \item Front-to-back ratio $\leq 0$\,dB.
    \item Azimuth beamwidth $\geq 360^\circ$.
\end{itemize}

% ── Dish ─────────────────────────────────────────────────────────
\subsection{Dish (Parabolic Reflector)}
\label{sec:dish}

\textbf{Registry name:} \texttt{Dish} \quad\textbar\quad \textbf{Has azimuth pattern:} Yes

\subsubsection{Default Parameters}
\begin{tabular}{lll}
\toprule
\textbf{Parameter} & \textbf{Default} & \textbf{Description} \\
\midrule
\texttt{dish\_efficiency} & 0.60 & Aperture efficiency (0 to 1) \\
\texttt{feed\_edge\_taper\_db} & $-12.0$ & Feed illumination taper at dish edge (dB) \\
\texttt{ftb\_ratio\_db} & 30.0 & Front-to-back ratio in dB \\
\texttt{max\_gain\_dbi} & 30.0 & Peak gain in dBi \\
\texttt{az\_beamwidth\_deg} & 5.0 & Azimuth 3\,dB beamwidth \\
\texttt{el\_beamwidth\_deg} & 5.0 & Elevation 3\,dB beamwidth \\
\texttt{f\_min\_mhz} & 3000.0 & Lower band edge in MHz \\
\texttt{f\_max\_mhz} & 6000.0 & Upper band edge in MHz \\
\bottomrule
\end{tabular}

\subsubsection{Mathematical Model}
The Dish antenna uses the Airy pattern (Jinc-squared) with an optional feed taper envelope:
\begin{equation}
G(\theta) = G_{\text{peak}} \cdot \eta \cdot \left[\frac{2\,J_1(u)}{u}\right]^2 \cdot T(r)
\end{equation}
where:
\begin{equation}
u = 1.6163 \cdot \frac{\sin(\theta)}{\sin(\theta_{3\text{dB}}/2)}
\end{equation}
$\eta$ is the aperture efficiency, $J_1$ is the Bessel function of the first kind (from SciPy when available; otherwise a polynomial approximation is used), and $T(r) = \exp(-r^2 / 2\sigma^2)$ is a Gaussian taper envelope derived from \texttt{feed\_edge\_taper\_db}. The taper sigma is chosen so that the taper amplitude equals $10^{(\text{taper\_dB}/20)}$ at the dish edge ($r = 1$). When \texttt{feed\_edge\_taper\_db} is $0$\,dB, $T(r) = 1$ (uniform illumination). The back lobe is weighted by the front-to-back ratio and elevation.

\subsubsection{Validation Warnings}
\begin{itemize}
    \item Aperture efficiency outside the range $0.3$--$0.85$.
    \item Azimuth beamwidth $\geq 90^\circ$.
    \item Gain $< 15$\,dBi.
\end{itemize}

% ── Horn ─────────────────────────────────────────────────────────
\subsection{Horn (Pyramidal)}
\label{sec:horn}

\textbf{Registry name:} \texttt{Horn} \quad\textbar\quad \textbf{Has azimuth pattern:} Yes

\subsubsection{Default Parameters}
\begin{tabular}{lll}
\toprule
\textbf{Parameter} & \textbf{Default} & \textbf{Description} \\
\midrule
\texttt{ftb\_ratio\_db} & 25.0 & Front-to-back ratio in dB \\
\texttt{horn\_aperture\_wavelengths} & \texttt{None} & Physical aperture width in wavelengths (optional) \\
\texttt{max\_gain\_dbi} & 15.0 & Peak gain in dBi \\
\texttt{az\_beamwidth\_deg} & 30.0 & Azimuth 3\,dB beamwidth \\
\texttt{el\_beamwidth\_deg} & 30.0 & Elevation 3\,dB beamwidth \\
\texttt{f\_min\_mhz} & 4000.0 & Lower band edge in MHz \\
\texttt{f\_max\_mhz} & 8000.0 & Upper band edge in MHz \\
\bottomrule
\end{tabular}

When \texttt{horn\_aperture\_wavelengths} is set, beamwidths are derived from the aperture size and scale with frequency: $BW_H = 67^\circ / (a \cdot f/f_{ref})$, $BW_E = 51^\circ / (a \cdot f/f_{ref})$. When \texttt{None}, beamwidths are set manually.

\subsubsection{Mathematical Model}
The Horn antenna uses a standard gain horn model with separate E-plane and H-plane formulations:

\textbf{E-plane (elevation):}
\begin{equation}
G_E = \operatorname{sinc}^2(u_{\text{el}}) \quad\text{where}\quad u_{\text{el}} = 1.3916 \cdot \frac{\sin(\theta_{\text{el}})}{\sin(\theta_{\text{el},3\text{dB}}/2)}
\end{equation}

\textbf{H-plane (azimuth):}
\begin{equation}
G_H = \left[\frac{\cos(v)}{1 - (2v/\pi)^2}\right]^2 \quad\text{where}\quad v = 1.1891 \cdot \frac{\sin(\theta_{\text{az}})}{\sin(\theta_{\text{az},3\text{dB}}/2)}
\end{equation}

The back lobe is weighted by the front-to-back ratio and elevation.

\subsubsection{Validation Warnings}
\begin{itemize}
    \item Azimuth beamwidth $\geq 180^\circ$.
    \item Gain $< 5$\,dBi.
\end{itemize}

% ── Monopole ─────────────────────────────────────────────────────
\subsection{Monopole / Dipole}
\label{sec:monopole}

\textbf{Registry name:} \texttt{Monopole} \quad\textbar\quad \textbf{Has azimuth pattern:} No (omnidirectional in azimuth)

\subsubsection{Default Parameters}
\begin{tabular}{lll}
\toprule
\textbf{Parameter} & \textbf{Default} & \textbf{Description} \\
\midrule
\texttt{element\_length\_wavelengths} & 0.25 & Element length in wavelengths \\
\texttt{monopole\_physical\_length\_m} & \texttt{None} & Physical element length in metres (optional) \\
\texttt{max\_gain\_dbi} & 2.15 & Peak gain in dBi \\
\texttt{f\_min\_mhz} & 100.0 & Lower band edge in MHz \\
\texttt{f\_max\_mhz} & 500.0 & Upper band edge in MHz \\
\bottomrule
\end{tabular}

When \texttt{monopole\_physical\_length\_m} is set, the electrical length scales with frequency: $L(f) = \text{physical\_length} / (c/f)$. This models a fixed physical element whose pattern evolves across the band. When \texttt{None}, every frequency slice uses the same \texttt{element\_length\_wavelengths}.

\subsubsection{Mathematical Model}
For a quarter-wave monopole ($L = 0.25\lambda$):
\begin{equation}
G(\theta_{\text{el}}) = G_{\text{peak}} \cdot \left[\frac{\cos\!\bigl(\frac{\pi}{2} \sin(\theta_{\text{el}})\bigr)}{\cos(\theta_{\text{el}})}\right]^2
\end{equation}

For a general dipole of length $L$ wavelengths:
\begin{equation}
G(\theta_{\text{el}}) = G_{\text{peak}} \cdot \frac{\bigl[\cos\!\bigl(\frac{kL}{2}\sin(\theta_{\text{el}})\bigr) - \cos\!\bigl(\frac{kL}{2}\bigr)\bigr]^2}{\bigl[1 - \cos\!\bigl(\frac{kL}{2}\bigr)\bigr]^2 \cdot \cos^2(\theta_{\text{el}})}
\end{equation}
where $kL = 2\pi L$.

\subsubsection{Validation Warnings}
\begin{itemize}
    \item Element length $> 1.5\lambda$ (multi-lobe patterns expected).
    \item Azimuth beamwidth $\neq 360^\circ$.
\end{itemize}

% ── Array ────────────────────────────────────────────────────────
\newpage
\subsection{Phased Array}
\label{sec:array}

\textbf{Registry name:} \texttt{Array} \quad\textbar\quad \textbf{Has azimuth pattern:} Yes

\subsubsection{Default Parameters}
\begin{tabular}{lll}
\toprule
\textbf{Parameter} & \textbf{Default} & \textbf{Description} \\
\midrule
\texttt{array\_geometry} & \texttt{linear} & Array layout: linear, planar, or circular \\
\texttt{array\_n\_elements\_x} & 8 & Number of elements along the X axis \\
\texttt{array\_n\_elements\_y} & 1 & Number of elements along the Y axis (planar) \\
\texttt{array\_spacing\_x\_lambda} & 0.5 & X-axis element spacing in wavelengths \\
\texttt{array\_spacing\_y\_lambda} & 0.5 & Y-axis element spacing in wavelengths \\
\texttt{array\_steer\_az\_deg} & 0.0 & Beam steering azimuth angle \\
\texttt{array\_steer\_el\_deg} & 0.0 & Beam steering elevation angle \\
\texttt{array\_weighting} & \texttt{uniform} & Amplitude taper: uniform or Taylor \\
\texttt{array\_taper\_sll\_db} & $-25.0$ & Taylor taper sidelobe level in dB \\
\texttt{array\_mutual\_coupling} & \texttt{False} & Enable mutual coupling model \\
\texttt{array\_element\_pattern} & \texttt{gaussian} & Element pattern type: gaussian, cosine, or patch \\
\texttt{array\_rms\_phase\_error\_deg} & 0.0 & RMS phase error per element (degrees) \\
\texttt{array\_rms\_amplitude\_error\_db} & 0.0 & RMS amplitude error per element (dB) \\
\texttt{array\_circular\_orientation} & \texttt{parallel} & Circular array element pointing: parallel or radial \\
\texttt{ftb\_ratio\_db} & 20.0 & Front-to-back ratio in dB \\
\texttt{max\_gain\_dbi} & 18.0 & Peak gain in dBi \\
\texttt{az\_beamwidth\_deg} & 30.0 & Azimuth 3\,dB beamwidth \\
\texttt{el\_beamwidth\_deg} & 30.0 & Elevation 3\,dB beamwidth \\
\texttt{f\_min\_mhz} & 2000.0 & Lower band edge in MHz \\
\texttt{f\_max\_mhz} & 4000.0 & Upper band edge in MHz \\
\bottomrule
\end{tabular}

\subsubsection{Mathematical Model}
The Phased Array uses the Pattern Multiplication Theorem:
\begin{equation}
G(\theta, \phi) = G_{\text{element}}(\theta, \phi) \cdot |AF(\theta, \phi)|^2
\end{equation}

\textbf{Element Pattern:} Three models are available via \texttt{array\_element\_pattern}:
\begin{itemize}
    \item \texttt{gaussian} (default): $G_{elem} = 0.5^{r^2}$ where $r^2 = (\theta_{az}/\theta_{hw})^2 + (\theta_{el}/\theta_{hw})^2$
    \item \texttt{cosine}: $G_{elem} = [\cos(\pi/2 \cdot \theta/\theta_{3dB})]^2$ per plane
    \item \texttt{patch}: $G_{elem} = \cos^{1.5}(\theta_{az}) \cdot \cos^{1.5}(\theta_{el})$ (empirical microstrip patch model)
\end{itemize}

\textbf{Linear Array Factor:}
\begin{equation}
AF(\theta) = \sum_{n=0}^{N-1} w_n \cdot e^{\,j\,n\,(k\,d\,\sin\theta + \beta)}
\end{equation}
where $w_n$ are the element weights (uniform or Taylor taper), $k = 2\pi/\lambda$, $d$ is the element spacing, and $\beta$ is the progressive phase shift for beam steering.

\textbf{Planar Array:} Formed as the product of two orthogonal linear array factors (X and Y axes).

\textbf{Circular Array:} Elements distributed uniformly on a circle; the array factor accounts for each element's position in polar coordinates. Element orientation is controlled by \texttt{array\_circular\_orientation}: \texttt{parallel} (all elements point broadside) or \texttt{radial} (each element points outward from centre, with gain weighted by $\cos(\phi - \phi_n)$).

\textbf{Mutual Coupling:} When enabled, a complex $N \times N$ impedance matrix is constructed using an exponential-decay model: $Z_{mn} = c_0 \cdot \exp(-\alpha \cdot |m-n| \cdot d/\lambda)$ for $m \neq n$, $Z_{mm} = 1$. The coupled weights are obtained by solving $Z \cdot w_{coupled} = w_{ideal}$ via Gaussian elimination with partial pivoting.

\textbf{Element Errors:} When \texttt{array\_rms\_amplitude\_error\_db} or \texttt{array\_rms\_phase\_error\_deg} are non-zero, Gaussian-distributed random perturbations are applied to each element's weight and phase offset using a deterministic seed for repeatability.

\subsubsection{Validation Warnings}
\begin{itemize}
    \item Element spacing $> 0.5\lambda$ in either axis (grating lobes may appear).
\end{itemize}

% ===================================================================
%  PART IV — PHYSICS FEATURES
% ===================================================================
\newpage
\section{Physics Features}
\label{sec:features}

AntennaForge provides nine configurable physics features that can be independently enabled or disabled. Each feature is controlled by a sub-dictionary within the \texttt{features} key of the configuration.

\subsection{Frequency-Dependent Beamwidth (Azimuth)}
\label{sec:feat_bw_az}

\textbf{Config key:} \texttt{features.freq\_dependent\_bw\_az}\\
\textbf{Default:} Enabled\\
\textbf{GUI control:} Switch (toggle on/off)

Causes the azimuth beamwidth to narrow as frequency increases, modelling the physical behaviour of real antennas. Three decay modes are available:

\subsubsection{Exponent Mode}
\begin{equation}
\mathrm{BW}(f) = \mathrm{BW}_{\text{base}} \cdot \left(\frac{f_{\text{ref}}}{f}\right)^{\alpha}
\end{equation}
where $\alpha$ is the scaling exponent (default: 0.8). For LPDA antennas, an additional damping factor (LPDA variation factor, default: 0.4) is applied to moderate the frequency dependence, reflecting the LPDA's broadband nature.

\begin{tabular}{lll}
\toprule
\textbf{Parameter} & \textbf{Default} & \textbf{Description} \\
\midrule
\texttt{scaling\_exponent} & 0.8 & Power-law exponent $\alpha$ \\
\texttt{lpda\_variation\_factor} & 0.4 & LPDA damping factor \\
\bottomrule
\end{tabular}

\subsubsection{Percentage Mode}
\begin{equation}
\mathrm{BW}(f) = \mathrm{BW}_{\text{base}} \cdot \left(1 - \frac{P}{100}\right)^{\log_2(f/f_{\text{ref}})}
\end{equation}
where $P$ is the percentage reduction per octave (default: 15\%).

\begin{tabular}{lll}
\toprule
\textbf{Parameter} & \textbf{Default} & \textbf{Description} \\
\midrule
\texttt{decay\_pct\_per\_octave} & 15.0 & Percent BW reduction per octave \\
\bottomrule
\end{tabular}

\subsubsection{Three-Point Mode}
A quadratic (Lagrange) interpolation through three user-specified beamwidth values at the start, middle, and end of the frequency band.

\begin{tabular}{lll}
\toprule
\textbf{Parameter} & \textbf{Default} & \textbf{Description} \\
\midrule
\texttt{three\_point\_start\_deg} & None & BW at $f_{\min}$ \\
\texttt{three\_point\_mid\_deg} & None & BW at $f_{\text{mid}}$ \\
\texttt{three\_point\_end\_deg} & None & BW at $f_{\max}$ \\
\bottomrule
\end{tabular}

\subsubsection{Beamwidth Clamps}
Regardless of the decay mode, the computed beamwidth is clamped between configurable minimum and maximum values:

\begin{tabular}{lll}
\toprule
\textbf{Parameter} & \textbf{Default} & \textbf{Description} \\
\midrule
\texttt{min\_bw\_deg} & 10.0 & Floor for azimuth beamwidth \\
\texttt{max\_bw\_deg} & 180.0 & Ceiling for azimuth beamwidth \\
\bottomrule
\end{tabular}

\subsection{Frequency-Dependent Beamwidth (Elevation)}
\label{sec:feat_bw_el}

\textbf{Config key:} \texttt{features.freq\_dependent\_bw\_el}\\
\textbf{Default:} Enabled

Identical to the azimuth version (Section~\ref{sec:feat_bw_az}) but applied independently to the elevation plane. Default elevation clamps are $10^\circ$--$90^\circ$.

\subsection{Sidelobes (Taylor)}
\label{sec:feat_sidelobes}

\textbf{Config key:} \texttt{features.sidelobes}\\
\textbf{Default:} Enabled\\
\textbf{GUI control:} Switch (toggle on/off)

Adds a realistic Taylor sidelobe structure to the pattern. The sidelobe envelope is computed as a function of off-axis angle, with configurable first sidelobe level, decay rate, and number of sidelobe pairs.

\begin{tabular}{lll}
\toprule
\textbf{Parameter} & \textbf{Default} & \textbf{Description} \\
\midrule
\texttt{type} & \texttt{taylor} & Sidelobe model type \\
\texttt{first\_sidelobe\_db} & $-13.2$ & First sidelobe level relative to peak (dB) \\
\texttt{decay\_rate\_db} & 2.0 & Sidelobe rolloff rate (dB per lobe) \\
\texttt{n\_sidelobes} & 5 & Number of sidelobe pairs per cut \\
\bottomrule
\end{tabular}

\subsection{Pattern Breakup}
\label{sec:feat_breakup}

\textbf{Config key:} \texttt{features.pattern\_breakup}\\
\textbf{Default:} Enabled\\
\textbf{GUI control:} Switch (toggle on/off)

Adds deterministic ripple at high off-axis angles to simulate the pattern irregularities observed in real antenna measurements.

\begin{tabular}{lll}
\toprule
\textbf{Parameter} & \textbf{Default} & \textbf{Description} \\
\midrule
\texttt{onset\_angle\_deg} & 90.0 & Angle from boresight where breakup begins \\
\texttt{ripple\_amplitude\_db} & 4.0 & Peak-to-peak ripple amplitude in dB \\
\texttt{ripple\_density} & 3.0 & Spatial frequency of the ripple \\
\bottomrule
\end{tabular}

\newpage
\subsection{Ground Reflection}
\label{sec:feat_ground}

\textbf{Config key:} \texttt{features.ground\_reflection}\\
\textbf{Default:} Disabled\\
\textbf{GUI control:} Switch (toggle on/off)

Models the 2-ray ground plane reflection interference pattern that occurs when an antenna is mounted above a reflective surface.

\begin{tabular}{lll}
\toprule
\textbf{Parameter} & \textbf{Default} & \textbf{Description} \\
\midrule
\texttt{height\_wavelengths} & 1.0 & Antenna height above ground in wavelengths \\
\texttt{reflection\_coeff} & 0.7 & Ground reflection coefficient (0 to 1) \\
\texttt{ground\_type} & \texttt{good\_soil} & Ground material type \\
\bottomrule
\end{tabular}

\subsection{Cross-Polarisation}
\label{sec:feat_xpol}

\textbf{Config key:} \texttt{features.cross\_pol}\\
\textbf{Default:} Disabled\\
\textbf{GUI control:} Switch (toggle on/off)

When enabled, generates additional cross-polarisation pattern CSV files alongside co-polarisation files. The cross-pol gain is modelled as a function of angular offset from boresight.

\begin{tabular}{lll}
\toprule
\textbf{Parameter} & \textbf{Default} & \textbf{Description} \\
\midrule
\texttt{boresight\_isolation\_db} & $-25.0$ & Cross-pol isolation at boresight \\
\texttt{peak\_angle\_deg} & 45.0 & Angle of peak cross-pol \\
\texttt{max\_cross\_pol\_db} & $-15.0$ & Maximum cross-pol level \\
\bottomrule
\end{tabular}

\subsection{VSWR Rolloff}
\label{sec:feat_vswr}

\textbf{Config key:} \texttt{features.vswr\_rolloff}\\
\textbf{Default:} Enabled\\
\textbf{GUI control:} Switch (toggle on/off)

Applies gain reduction at the band edges to model the impedance mismatch that occurs near the frequency limits of an antenna's operational bandwidth.

\begin{tabular}{lll}
\toprule
\textbf{Parameter} & \textbf{Default} & \textbf{Description} \\
\midrule
\texttt{rolloff\_band\_fraction} & 0.1 & Fraction of band affected at each edge \\
\texttt{max\_rolloff\_db} & 3.0 & Maximum gain reduction at band edge \\
\texttt{rolloff\_shape} & \texttt{cosine} & Rolloff curve shape \\
\bottomrule
\end{tabular}

\newpage
\subsection{Asymmetry / Squint}
\label{sec:feat_asymmetry}

\textbf{Config key:} \texttt{features.asymmetry}\\
\textbf{Default:} Disabled\\
\textbf{GUI control:} Switch (toggle on/off)

Introduces pointing error and deterministic left/right asymmetry to simulate real-world manufacturing tolerances and mounting imperfections.

\begin{tabular}{lll}
\toprule
\textbf{Parameter} & \textbf{Default} & \textbf{Description} \\
\midrule
\texttt{az\_squint\_deg} & 0.0 & Azimuth pointing offset in degrees \\
\texttt{el\_tilt\_deg} & 0.0 & Elevation tilt offset in degrees \\
\texttt{random\_asymmetry\_db} & 1.0 & Peak random perturbation in dB \\
\bottomrule
\end{tabular}

\subsection{Gain-Beamwidth Coupling}
\label{sec:feat_gbw}

\textbf{Config key:} \texttt{features.gain\_bw\_coupling}\\
\textbf{Default:} Disabled\\
\textbf{GUI control:} Switch (toggle on/off) with mode radio buttons

Links frequency-dependent gain and beamwidth behaviour so that both vary consistently across the band. Three coupling modes are available:

\begin{description}[style=nextline, leftmargin=2em]
    \item[\textbf{Gain drives beamwidth}] The gain rolloff with frequency drives a corresponding beamwidth change via the directivity-beamwidth relationship.
    \item[\textbf{Beamwidth drives gain}] The beamwidth scaling with frequency determines the gain variation.
    \item[\textbf{Independent curves}] Both gain and beamwidth vary with frequency according to their own independent parameters (default).
\end{description}

\begin{tabular}{lll}
\toprule
\textbf{Parameter} & \textbf{Default} & \textbf{Description} \\
\midrule
\texttt{mode} & \texttt{independent} & Coupling mode \\
\texttt{gain\_rolloff\_db\_per\_octave} & 1.5 & Gain decrease per octave below $f_{\max}$ \\
\bottomrule
\end{tabular}

In the GUI, the Gain-Beamwidth Coupling card appears for all directional antenna types (LPDA, Panel, Dish, Horn, Array). It is automatically enabled when LPDA is selected.

\subsection{Front-Fade Steepness (Sigmoid $k$)}
\label{sec:sigmoid}

\textbf{Config key:} \texttt{sigmoid\_k} (top-level)\\
\textbf{Default:} 40.0\\
\textbf{GUI control:} Slider + Entry (range: 2--60, step: 1)

Controls the sharpness of the smooth sigmoid transition between the front hemisphere and back lobe. A higher $k$ value produces a sharper transition; a lower value creates a more gradual blend. The sigmoid function is:
\begin{equation}
w(\cos\theta_{\text{az}}) = \frac{1}{1 + e^{-k \cdot \cos\theta_{\text{az}}}}
\end{equation}

% ===================================================================
%  PART V — GUI REFERENCE
% ===================================================================
\newpage
\section{Graphical User Interface (GUI)}
\label{sec:gui}

The AntennaForge GUI is built with \textbf{customtkinter} (CTk). The title bar reads ``AntennaForge v1.0''.

\subsection{Layout}
The application uses a \textbf{sidebar + stacked content area} layout:
\begin{itemize}
    \item \textbf{Left sidebar} — Navigation buttons for all nine pages. Each button displays an icon and label. The active page is highlighted in blue.
    \item \textbf{Main content area} — Displays the currently selected page. Pages are lazy-loaded on first visit and cached for performance.
    \item \textbf{Status bar} — A bottom bar spanning both columns, showing messages and a progress bar.
\end{itemize}

\subsection{Sidebar Navigation}
\begin{longtable}{cllL{6cm}}
\toprule
\textbf{Icon} & \textbf{Label} & \textbf{Page Key} & \textbf{Description} \\
\midrule
\endhead
Home & Home & \texttt{home} & Dashboard with quick-start actions \\
Antenna & Generate & \texttt{generate} & Pattern generation (single \& batch) \\
Chart & Visualize & \texttt{visualize} & Plot patterns (heatmap, cuts, polar, 3D) \\
Magnify & Analysis & \texttt{analysis} & Pattern analysis and comparison \\
Gear & Operations & \texttt{operations} & Arithmetic, rotation, interpolation \\
Ruler & RF Calcs & \texttt{rf\_calc} & EIRP, FSPL, link budget \\
Document & Reports & \texttt{reports} & PDF and \LaTeX\ reports \\
Help & How To & \texttt{howto} & Built-in user guide \\
Wrench & Settings & \texttt{settings} & Theme, fonts, colours, config I/O \\
\bottomrule
\end{longtable}

\subsection{Theme System}
\label{sec:theme}

The theme engine (\texttt{gui/theme.py}) defines a 27-colour palette, each stored as a (light mode, dark mode) tuple. Colours are grouped into six categories: \textbf{Sidebar} (background, hover, active, text), \textbf{Main} (background, surface, card), \textbf{Text} (primary, secondary, muted), \textbf{Accent} (blue highlights), \textbf{Status} (success, warning, error, info), and \textbf{Inputs} (entry background, border, focus). All colours can be customised per-channel via the Settings page (see Section~\ref{sec:colour_custom}).

\subsubsection{Font Definitions}
Fonts are built from three base parameters: \texttt{family} (default: \texttt{Segoe UI}), \texttt{mono} (default: \texttt{Consolas}), and \texttt{base} size (default: 12\,pt). All other font sizes (headings, subheadings, small text, etc.) are derived from the base size. Fonts can be changed via the Settings page.


\subsection{User Preferences}
\label{sec:prefs}

User preferences are persisted in \texttt{antennaforge\_prefs.json} at the project root. The file stores:
\begin{lstlisting}[language={},caption={Preferences file structure}]
{
    "theme": "dark",
    "colors": { "<key>": ["#light", "#dark"], ... },
    "fonts": { "family": "Segoe UI", "mono": "Consolas", "base": 12 }
}
\end{lstlisting}

% ===================================================================
%  PART VI — GUI PAGES IN DETAIL
% ===================================================================
\newpage
\section{GUI Pages — Complete Field Reference}
\label{sec:pages}

\subsection{Home Page}
\label{sec:page_home}

The Home page serves as the dashboard landing screen with quick-start navigation cards.

\subsubsection{Quick-Action Cards}
Four cards in a row, each with an icon, title, description, and an ``Open $\rightarrow$'' button:

\begin{tabular}{clll}
\toprule
\textbf{Icon} & \textbf{Title} & \textbf{Target} & \textbf{Description} \\
\midrule
"Antenna" & Generate Patterns & Generate & Create antenna pattern CSVs \\
"Chart" & Visualize Patterns & Visualize & Heatmaps, cuts, polar \& 3D \\
"Magnify" & Analyze Patterns & Analysis & Beamwidth, sidelobes, symmetry \\
"Document" & Generate Reports & Reports & PDF and \LaTeX\ reports \\
\bottomrule
\end{tabular}

\subsubsection{Tools Row}
Three secondary cards for Pattern Ops, RF Calculations, and Settings.

\subsubsection{Capabilities Card}
Lists eight bullet-point capabilities of the application.

\subsubsection{Dependency Status Card}
Dynamically checks installed packages and displays a status indicator ($\checkmark$ or $\times$) for each optional dependency. Shows the pip install command for any missing packages.

% ── GENERATE PAGE ────────────────────────────────────────────────
\newpage
\subsection{Generate Page}
\label{sec:page_generate}

The Generate page is the primary antenna pattern creation interface. It contains multiple collapsible cards with context-sensitive widget visibility based on the selected antenna type.

\subsubsection{Dynamic Visibility Rules}
\label{sec:visibility}

Not all controls are visible for all antenna types. The following table shows which widget groups appear for each type:

\begin{longtable}{L{4cm} C{0.7cm} C{0.7cm} C{0.7cm} C{0.7cm} C{0.7cm} C{1.0cm} C{0.7cm}}
\toprule
\textbf{Widget Group} & \rotatebox{90}{\textbf{LPDA}} & \rotatebox{90}{\textbf{Omni}} & \rotatebox{90}{\textbf{Panel}} & \rotatebox{90}{\textbf{Dish}} & \rotatebox{90}{\textbf{Horn}} & \rotatebox{90}{\textbf{Monopole}} & \rotatebox{90}{\textbf{Array}} \\
\midrule
\endhead
Front-to-Back Ratio & $\checkmark$ & & $\checkmark$ & $\checkmark$ & $\checkmark$ & & $\checkmark$ \\
Azimuth Beamwidth & $\checkmark$ & & $\checkmark$ & $\checkmark$ & $\checkmark$ & & $\checkmark$ \\
Elevation Beamwidth & $\checkmark$ & $\checkmark$ & $\checkmark$ & $\checkmark$ & $\checkmark$ & & $\checkmark$ \\
Downtilt & & & $\checkmark$ & & & & \\
Aperture Efficiency & & & & $\checkmark$ & & & \\
Element Length & & & & & & $\checkmark$ & \\
Sidelobe Parameters & $\checkmark$ & $\checkmark$ & & & & & \\
Front-Fade Steepness & $\checkmark$ & & $\checkmark$ & $\checkmark$ & $\checkmark$ & & \\
Array Parameters & & & & & & & $\checkmark$ \\
Freq-dep BW (Az) & $\checkmark$ & & $\checkmark$ & $\checkmark$ & $\checkmark$ & & $\checkmark$ \\
Freq-dep BW (El) & $\checkmark$ & $\checkmark$ & $\checkmark$ & $\checkmark$ & $\checkmark$ & & $\checkmark$ \\
Pattern Breakup & $\checkmark$ & $\checkmark$ & $\checkmark$ & $\checkmark$ & $\checkmark$ & & $\checkmark$ \\
Ground Reflection & $\checkmark$ & $\checkmark$ & $\checkmark$ & $\checkmark$ & $\checkmark$ & $\checkmark$ & $\checkmark$ \\
Cross-Polarisation & $\checkmark$ & $\checkmark$ & $\checkmark$ & $\checkmark$ & $\checkmark$ & $\checkmark$ & $\checkmark$ \\
VSWR Rolloff & $\checkmark$ & $\checkmark$ & $\checkmark$ & $\checkmark$ & $\checkmark$ & $\checkmark$ & $\checkmark$ \\
Asymmetry / Squint & $\checkmark$ & & $\checkmark$ & $\checkmark$ & $\checkmark$ & & $\checkmark$ \\
Gain-BW Coupling & $\checkmark$ & & $\checkmark$ & $\checkmark$ & $\checkmark$ & & $\checkmark$ \\
\bottomrule
\end{longtable}

\subsubsection{Card: Core Parameters}

\begin{longtable}{L{3.2cm} L{2.5cm} L{1.5cm} L{6cm}}
\toprule
\textbf{Field} & \textbf{Widget Type} & \textbf{Default} & \textbf{Description} \\
\midrule
\endhead
Antenna Type & Dropdown & LPDA & Selects the antenna model. Options: LPDA, Omni, Panel, Dish, Horn, Monopole, Array. Changing this triggers dynamic visibility updates for all other fields. \\
Max Gain (dBi) & Text Entry & 7.0 & Peak boresight gain of the antenna in dBi. Valid range: $-30$ to $60$\,dBi. \\
Front-to-Back (dB) & Text Entry & 15.0 & Front-to-back ratio in dB. Controls the gain difference between the front and rear of the antenna. Valid range: $0$ to $80$\,dB. Hidden for Omni and Monopole. \\
Downtilt ($^\circ$) & Text Entry & 0.0 & Electrical downtilt angle in degrees. Positive values tilt the beam downward. Panel type only. \\
Aperture Efficiency & Text Entry & 0.60 & Ratio of effective to physical aperture area (0 to 1). Typical values: 0.55--0.65. Dish type only. \\
Element Length ($\lambda$) & Text Entry & 0.25 & Monopole/dipole element length in wavelengths. Common values: 0.25 (quarter-wave), 0.5 (half-wave). Monopole type only. \\
\bottomrule
\end{longtable}

\subsubsection{Card: Beamwidth}

\begin{longtable}{L{3.2cm} L{3cm} L{1.5cm} L{1.5cm} L{4cm}}
\toprule
\textbf{Field} & \textbf{Widget Type} & \textbf{Default} & \textbf{Range} & \textbf{Description} \\
\midrule
\endhead
Azimuth BW ($^\circ$) & Slider + Entry & 65 & 1--360 & The 3\,dB azimuth beamwidth at the reference frequency. The slider and entry field are bidirectionally synchronised — moving the slider updates the entry and typing in the entry moves the slider. \\
Elevation BW ($^\circ$) & Slider + Entry & 70 & 1--360 & The 3\,dB elevation beamwidth at the reference frequency. Bidirectionally synchronised slider and entry. \\
\bottomrule
\end{longtable}

\subsubsection{Card: Array Parameters}
\label{sec:array_params}
\textit{Visible only when Antenna Type is set to Array.}

\begin{longtable}{L{3.2cm} L{2.2cm} L{1.5cm} L{6.3cm}}
\toprule
\textbf{Field} & \textbf{Widget Type} & \textbf{Default} & \textbf{Description} \\
\midrule
\endhead
Geometry & Dropdown & linear & Array layout geometry. Options: \textbf{linear} (1D), \textbf{planar} (2D rectangular grid), \textbf{circular} (elements on a ring). \\
Elements X & Text Entry & 8 & Number of elements along the X axis. For circular arrays, this is the total number of elements. \\
Elements Y & Text Entry & 1 & Number of elements along the Y axis. Only meaningful for planar geometry. \\
Spacing X ($\lambda$) & Text Entry & 0.5 & Inter-element spacing along X in wavelengths. Values $> 0.5$ may produce grating lobes. \\
Spacing Y ($\lambda$) & Text Entry & 0.5 & Inter-element spacing along Y in wavelengths. Only meaningful for planar geometry. \\
Steer Az ($^\circ$) & Text Entry & 0.0 & Beam steering angle in azimuth. The array factor's progressive phase shift is computed from this angle. \\
Steer El ($^\circ$) & Text Entry & 0.0 & Beam steering angle in elevation. \\
Weighting & Dropdown & uniform & Element amplitude tapering. \textbf{uniform}: equal weights. \textbf{taylor}: Taylor taper for reduced sidelobes. \\
Taper SLL (dB) & Text Entry & $-25.0$ & Desired sidelobe level for Taylor taper weighting. Only meaningful when Weighting is set to ``taylor''. \\
Mutual Coupling & Switch (toggle) & OFF & When enabled, a complex impedance matrix is computed based on inter-element distances and applied as weight corrections. Increases computation time. \\
\bottomrule
\end{longtable}

\subsubsection{Card: Frequency}

\begin{longtable}{L{3.5cm} L{2.5cm} L{1.5cm} L{5.7cm}}
\toprule
\textbf{Field} & \textbf{Widget Type} & \textbf{Default} & \textbf{Description} \\
\midrule
\endhead
F min (MHz) & Text Entry & 30.0 & Lower edge of the frequency band in MHz. Valid range: 0.001 to $10^6$\,MHz. \\
F max (MHz) & Text Entry & 300.0 & Upper edge of the frequency band in MHz. Must be $\geq$ F min. \\
Number of Slices & Text Entry & 10 & Number of frequency points to generate. Each slice produces one CSV file. Valid range: 1 to 10{,}000. \\
Freq Slice Spacing & Dropdown & linear & Distribution of frequency points across the band. \textbf{linear}: equally spaced. \textbf{log}: logarithmically spaced (more points at lower frequencies). \\
Ref Frequency (MHz) & Text Entry & (blank) & Reference frequency for beamwidth scaling. If left blank, the geometric mean $\sqrt{f_{\min} \cdot f_{\max}}$ is used automatically. \\
\bottomrule
\end{longtable}

\newpage
\subsubsection{Card: Angular Resolution}

\begin{longtable}{L{3.2cm} L{2.5cm} L{1.5cm} L{6cm}}
\toprule
\textbf{Field} & \textbf{Widget Type} & \textbf{Default} & \textbf{Description} \\
\midrule
\endhead
Azimuth Step ($^\circ$) & Dropdown & 1.0 & Angular step size for the azimuth grid. Options: 0.1, 0.5, 1.0, 2.0, 5.0. Smaller steps produce higher-resolution patterns but increase file size and computation time. \\
Elevation Step ($^\circ$) & Dropdown & 1.0 & Angular step size for the elevation grid. Same options as azimuth. \\
\bottomrule
\end{longtable}

\subsubsection{Card: Polarization}

\begin{longtable}{L{3.2cm} L{2.5cm} L{1.5cm} L{6cm}}
\toprule
\textbf{Field} & \textbf{Widget Type} & \textbf{Default} & \textbf{Description} \\
\midrule
\endhead
Polarization & Dropdown & vertical & Polarisation plane of the antenna. Options: \textbf{vertical}, \textbf{horizontal}, \textbf{custom}. Selecting ``custom'' reveals the Pol Angle field. \\
Pol Angle ($^\circ$) & Text Entry & 0.0 & Custom polarisation angle in degrees. Only visible when Polarization is set to ``custom''. $0^\circ$ = vertical, $90^\circ$ = horizontal. \\
Noise Mode & Radio Buttons & uniform & Noise distribution type: \textbf{uniform} ($\pm$ range), \textbf{gaussian} ($\sigma$ in dB), or \textbf{quantization} (step size in dB). \\
Noise Value (dB) & Slider + Entry & 0 & Noise amplitude applied to the pattern. Interpretation depends on mode: range for uniform, sigma for Gaussian, step size for quantization. Range: 0--5\,dB, step: 0.1\,dB. \\
\bottomrule
\end{longtable}

\subsubsection{Card: Physics Features}
\label{sec:gui_features}

This card contains boolean toggle switches for each physics feature. When toggled on, the corresponding feature is applied during pattern generation.

\begin{longtable}{L{4cm} L{1.5cm} L{7.7cm}}
\toprule
\textbf{Switch Label} & \textbf{Default} & \textbf{Description} \\
\midrule
\endhead
Freq-dependent BW (Az) & ON & Scales the azimuth beamwidth with frequency. See Section~\ref{sec:feat_bw_az}. \\
Freq-dependent BW (El) & ON & Scales the elevation beamwidth with frequency. See Section~\ref{sec:feat_bw_el}. \\
Sidelobes (Taylor) & ON & Adds Taylor sidelobe structure. See Section~\ref{sec:feat_sidelobes}. \\
Pattern Breakup & OFF & Adds deterministic ripple at high off-axis angles. See Section~\ref{sec:feat_breakup}. \\
Ground Reflection & OFF & Enables 2-ray ground reflection model. See Section~\ref{sec:feat_ground}. \\
Cross-Polarization & OFF & Generates cross-pol CSV files. See Section~\ref{sec:feat_xpol}. \\
VSWR Rolloff & ON & Applies band-edge gain reduction. See Section~\ref{sec:feat_vswr}. \\
Asymmetry / Squint & OFF & Introduces pointing error and asymmetry. See Section~\ref{sec:feat_asymmetry}. \\
\bottomrule
\end{longtable}

\subsubsection{Card: Gain-Beamwidth Coupling}
\textit{Visible for all directional types (LPDA, Panel, Dish, Horn, Array). Automatically enabled for LPDA.}

\begin{longtable}{L{3.5cm} L{2.5cm} L{1.5cm} L{5.7cm}}
\toprule
\textbf{Field} & \textbf{Widget Type} & \textbf{Default} & \textbf{Description} \\
\midrule
\endhead
Enable Coupling & Switch (toggle) & OFF & When enabled, gain and beamwidth vary together across frequency. Auto-enabled for LPDA. \\
Mode & Radio Buttons & Independent & Three modes: \textbf{Gain drives beamwidth}, \textbf{Beamwidth drives gain}, \textbf{Independent curves}. Visible only when coupling is enabled. \\
Gain Rolloff (dB/octave) & Text Entry & 1.5 & How much peak gain decreases per octave below $f_{\max}$. Visible when mode is ``Gain drives beamwidth'' or ``Independent curves''. \\
\bottomrule
\end{longtable}

\subsubsection{Card: Sidelobe Parameters}
\textit{Visible for LPDA and Omni types only.}

\begin{longtable}{L{3.5cm} L{2cm} L{1.5cm} L{6.2cm}}
\toprule
\textbf{Field} & \textbf{Widget Type} & \textbf{Default} & \textbf{Description} \\
\midrule
\endhead
First SLL (dB) & Text Entry & $-20.0$ & First sidelobe level relative to the main beam peak, in dB. More negative values mean deeper sidelobes. \\
Decay Rate (dB/lobe) & Text Entry & 5.0 & Rate at which successive sidelobes decrease in amplitude, in dB per lobe. \\
Number of Sidelobes & Text Entry & 5 & Number of sidelobe pairs to generate on each side of the main beam in each principal plane. \\
\bottomrule
\end{longtable}

\subsubsection{Card: Front-Fade Steepness}

\begin{longtable}{L{3.5cm} L{3cm} L{1cm} L{1.5cm} L{4cm}}
\toprule
\textbf{Field} & \textbf{Widget Type} & \textbf{Def.} & \textbf{Range} & \textbf{Description} \\
\midrule
\endhead
Front-Fade Steepness ($k$) & Slider + Entry & 40 & 2--60 & Controls the sigmoid transition sharpness between front hemisphere and back lobe. Higher values = sharper transition. \\
\bottomrule
\end{longtable}

\newpage
\subsubsection{Card: Beamwidth vs Frequency}
\label{sec:bw_vs_freq_card}

This card configures how beamwidth changes across frequency for both azimuth and elevation planes independently.

\paragraph{Decay Mode Selection}

\begin{longtable}{L{3.2cm} L{2.5cm} L{1.5cm} L{6cm}}
\toprule
\textbf{Field} & \textbf{Widget Type} & \textbf{Default} & \textbf{Description} \\
\midrule
\endhead
Az Decay Mode & Dropdown & Exponent & Beamwidth scaling mode for azimuth. Options: Exponent, Percentage, Three-Point. \\
El Decay Mode & Dropdown & Exponent & Beamwidth scaling mode for elevation. Same options. \\
\bottomrule
\end{longtable}

\paragraph{Exponent Sub-Parameters}
\textit{Visible when the corresponding decay mode is ``Exponent''.}

\begin{longtable}{L{3.8cm} L{2cm} L{1cm} L{6.2cm}}
\toprule
\textbf{Field} & \textbf{Widget Type} & \textbf{Def.} & \textbf{Description} \\
\midrule
\endhead
Az Scaling Exponent & Text Entry & 0.8 & Power-law exponent for azimuth BW scaling: $\mathrm{BW} = \mathrm{BW}_0 \cdot (f_{\text{ref}}/f)^\alpha$. \\
El Scaling Exponent & Text Entry & 0.8 & Same for elevation. \\
Az LPDA Variation Factor & Text Entry & 0.3 & Damping factor for LPDA-type antennas (visible only when LPDA + Exponent). \\
El LPDA Variation Factor & Text Entry & 0.3 & Same for elevation. \\
\bottomrule
\end{longtable}

\paragraph{Percentage Sub-Parameters}
\textit{Visible when the corresponding decay mode is ``Percentage''.}

\begin{longtable}{L{3.8cm} L{2cm} L{1cm} L{6.2cm}}
\toprule
\textbf{Field} & \textbf{Widget Type} & \textbf{Def.} & \textbf{Description} \\
\midrule
\endhead
Az \% per Octave & Text Entry & 15.0 & Percentage beamwidth reduction per frequency doubling. \\
El \% per Octave & Text Entry & 15.0 & Same for elevation. \\
\bottomrule
\end{longtable}

\paragraph{Three-Point Sub-Parameters}
\textit{Visible when the corresponding decay mode is ``Three-Point''.}

\begin{longtable}{L{3.8cm} L{2cm} L{1cm} L{6.2cm}}
\toprule
\textbf{Field} & \textbf{Widget Type} & \textbf{Def.} & \textbf{Description} \\
\midrule
\endhead
Az BW @ $f_{\text{start}}$ (deg) & Text Entry & 90 & Azimuth beamwidth at the lowest frequency. \\
Az BW @ $f_{\text{mid}}$ (deg) & Text Entry & 65 & Azimuth beamwidth at band centre. \\
Az BW @ $f_{\text{end}}$ (deg) & Text Entry & 45 & Azimuth beamwidth at the highest frequency. \\
El BW @ $f_{\text{start}}$ (deg) & Text Entry & 90 & Elevation beamwidth at $f_{\min}$. \\
El BW @ $f_{\text{mid}}$ (deg) & Text Entry & 70 & Elevation beamwidth at $f_{\text{mid}}$. \\
El BW @ $f_{\text{end}}$ (deg) & Text Entry & 50 & Elevation beamwidth at $f_{\max}$. \\
\bottomrule
\end{longtable}

\paragraph{Beamwidth Clamps}
\textit{Always visible regardless of decay mode.}

\begin{longtable}{L{3.2cm} L{2cm} L{1.2cm} L{6.8cm}}
\toprule
\textbf{Field} & \textbf{Widget Type} & \textbf{Default} & \textbf{Description} \\
\midrule
\endhead
Az Min BW (deg) & Text Entry & 10 & Minimum allowed azimuth beamwidth after frequency scaling. \\
Az Max BW (deg) & Text Entry & 180 & Maximum allowed azimuth beamwidth. \\
El Min BW (deg) & Text Entry & 10 & Minimum allowed elevation beamwidth. \\
El Max BW (deg) & Text Entry & 90 & Maximum allowed elevation beamwidth. \\
\bottomrule
\end{longtable}

\subsubsection{Card: Output}

\begin{longtable}{L{3.2cm} L{2.5cm} L{3.5cm} L{4cm}}
\toprule
\textbf{Field} & \textbf{Widget Type} & \textbf{Default} & \textbf{Description} \\
\midrule
\endhead
Output Directory & FilePicker (directory) & \texttt{../antenna\_patterns} & The directory where generated CSV files will be saved. Every run creates a timestamped subdirectory (\texttt{yyyymmddhhmm}), including the first run. \\
\bottomrule
\end{longtable}

\subsubsection{Action Buttons}

\begin{longtable}{L{3.5cm} L{2.5cm} L{7.2cm}}
\toprule
\textbf{Button} & \textbf{Style} & \textbf{Action} \\
\midrule
\endhead
Generate Patterns & Success (green) & Builds a configuration dictionary from all fields, validates it, and runs pattern generation in a background thread. Progress is displayed in the result box and status bar. \\
Reset Defaults & Secondary (grey) & Resets all fields on the Generate page to their factory default values. \\
\bottomrule
\end{longtable}

\subsubsection{Result Box}
A read-only text area that displays the generation progress log, including each frequency slice as it is computed, the output file path, and any errors or warnings.

\subsubsection{Card: Batch Processing}
\label{sec:batch}

Below the single-generation controls, the Generate page includes a Batch Processing card for running multiple antenna configurations from a JSON recipe.

\begin{longtable}{L{3cm} L{3cm} L{7.2cm}}
\toprule
\textbf{Field} & \textbf{Widget Type} & \textbf{Description} \\
\midrule
\endhead
Recipe JSON & FilePicker & Select a batch recipe JSON file. \\
Output Folder & FilePicker (directory) & Base output directory \textemdash\ each job creates a subfolder here. Leave blank to use each job's own \texttt{output\_dir}. \\
Generate Images & Toggle Switch & When enabled, plot images (heatmaps, polar, 3D surface, cuts, trend plots) are generated for every batch job output. A note warns that image generation is slow for large recipes. \\
Generate Report & Toggle Switch & Enabled only when Generate Images is on. When active, a PDF report is generated for each completed job using the selected preset. \\
Report Preset & Radio Buttons & Visible when Generate Report is on. Selects the report section preset: Default, Full, Minimal, or Quick (see Section~\ref{sec:report-config}). \\
\bottomrule
\end{longtable}

\textbf{Buttons:}
\begin{itemize}
    \item ``Run Batch'' (success/green) \textemdash\ Loads and executes the batch recipe, generating patterns for each configuration entry. Runs in a background thread.
    \item ``Cancel'' (danger/red) \textemdash\ Requests cooperative cancellation of the running batch. The current job finishes, then the loop stops. Disabled when no batch is running.
    \item ``Dry Run'' (secondary) \textemdash\ Validates the recipe without generating any files. Reports whether all configurations are valid and what would be generated.
    \item ``Sample Recipe'' (secondary) \textemdash\ Opens a save-as dialog and writes a starter template recipe JSON file.
\end{itemize}

Cancellation is cooperative: the current job runs to completion, then the batch stops.

A batch recipe JSON file is an array of configuration dictionaries. Each entry inherits from \texttt{DEFAULT\_CONFIG} and antenna-specific defaults, with user overrides merged on top. Example structure:

\begin{lstlisting}[language={},caption={Batch recipe example}]
[
    {
        "antenna_type": "LPDA",
        "max_gain_dbi": 10.0,
        "f_min_mhz": 30.0,
        "f_max_mhz": 300.0,
        "n_slices": 10,
        "output_dir": "./batch_output/lpda"
    },
    {
        "antenna_type": "Dish",
        "max_gain_dbi": 35.0,
        "dish_efficiency": 0.65,
        "f_min_mhz": 4000.0,
        "f_max_mhz": 6000.0,
        "n_slices": 20,
        "output_dir": "./batch_output/dish"
    }
]
\end{lstlisting}

% ── VISUALIZE PAGE ───────────────────────────────────────────────
\subsection{Visualize Page}
\label{sec:page_visualize}

The Visualize page renders antenna patterns using embedded Matplotlib. It supports caching of generated plot images, navigation between multiple plots, and cooperative cancellation of long-running render tasks.

\subsubsection{Card: Input}

\begin{longtable}{L{3.2cm} L{3cm} L{7cm}}
\toprule
\textbf{Field} & \textbf{Widget Type} & \textbf{Description} \\
\midrule
\endhead
Pattern CSV & FilePicker (file) & Select a single CSV file for Heatmap, Cuts, Polar, or 3D Surface plots. \\
Pattern Directory & FilePicker (directory) & Select a folder of CSVs for multi-file plots (Gain vs Freq, BW vs Freq, Overlay Cuts) or for iterating single-file plots across all files. \\
\bottomrule
\end{longtable}

\subsubsection{Card: Plot Type}

\begin{longtable}{L{3.2cm} L{2.5cm} L{1.5cm} L{6cm}}
\toprule
\textbf{Field} & \textbf{Widget Type} & \textbf{Default} & \textbf{Description} \\
\midrule
\endhead
Type & Dropdown & Heatmap & The type of plot to generate. Options described below. \\
GIF Style & Dropdown & heatmap & Style for animated GIF. Options: \textbf{heatmap}, \textbf{polar}. Only visible when Type is ``Freq Sweep (GIF)''. \\
\bottomrule
\end{longtable}

\paragraph{Plot Type Options}
\begin{description}[style=nextline, leftmargin=2em]
    \item[\textbf{Heatmap}] A 2D colour-mapped plot of gain (dBi) across azimuth (x-axis) and elevation (y-axis). Requires a single CSV file. Colour scale uses the matplotlib \texttt{jet} colourmap.
    \item[\textbf{Cuts (Az+El)}] Two sub-plots showing the azimuth cut at $\mathrm{el}=0^\circ$ and the elevation cut at $\mathrm{az}=0^\circ$. Annotates the 3\,dB beamwidth. Requires a single CSV.
    \item[\textbf{Polar}] Polar-coordinate plots of the azimuth and elevation principal-plane cuts. Requires a single CSV.
    \item[\textbf{3D Surface}] A three-dimensional surface rendering of the full gain pattern. Requires a single CSV. Uses perspective projection with interactive viewpoint.
    \item[\textbf{Gain vs Freq}] Plots peak gain (dBi) versus frequency across all CSVs in the selected directory. Requires a directory.
    \item[\textbf{BW vs Freq}] Plots measured azimuth and elevation beamwidths versus frequency. Requires a directory.
    \item[\textbf{Overlay Cuts}] Overlays the azimuth cuts from multiple frequency slices on a single plot for comparison. Requires a directory.
    \item[\textbf{Freq Sweep (GIF)}] Generates an animated GIF cycling through all frequency slices. Requires a directory. The GIF Style dropdown selects between heatmap and polar frame styles.
\end{description}

\subsubsection{Action Buttons}

\begin{longtable}{L{3cm} L{2cm} L{8.2cm}}
\toprule
\textbf{Button} & \textbf{Style} & \textbf{Action} \\
\midrule
\endhead
Plot & Primary & Dispatches to the appropriate single-file, multi-file, directory, or animation handler based on the selected plot type and input source. \\
Cancel & Secondary & Sets a cooperative cancellation flag. Long-running batch plot operations check this flag between iterations. \\
Open Folder & Secondary & Opens the plot output directory in Windows Explorer. \\
\bottomrule
\end{longtable}

\subsubsection{Plot Display Area}
\begin{itemize}
    \item \textbf{Navigation bar:} ``$\triangleleft$ Prev'' and ``Next $\triangleright$'' buttons with a counter label (e.g., ``1 / 5 --- filename.png'').
    \item \textbf{Image frame:} Displays the rendered plot as an embedded image, auto-scaled to fit the frame.
    \item \textbf{Placeholder:} Shows ``Select a file or directory and click Plot'' before any plot is generated.
\end{itemize}

\subsubsection{Result Box}
A read-only text area showing plot generation status and any errors.

% ── ANALYSIS PAGE ────────────────────────────────────────────────
\subsection{Analysis Page}
\label{sec:page_analysis}

The Analysis page provides six analysis modes. The \textbf{Parameter Fit} card is presented first (primary workflow), followed by \textbf{Batch Fitting}, single-pattern analysis, multi-pattern comparison, pattern-set comparison (percentage differences), and sidelobe mask compliance testing.

\subsubsection{Card: Single Pattern Analysis}

\begin{longtable}{L{3cm} L{2.5cm} L{7.7cm}}
\toprule
\textbf{Field} & \textbf{Widget Type} & \textbf{Description} \\
\midrule
\endhead
Pattern CSV & FilePicker (file) & The pattern CSV to analyse. \\
\bottomrule
\end{longtable}

\textbf{Button:} ``Analyze'' (primary) — Loads the CSV and computes:
\begin{itemize}
    \item Peak gain (maximum value in the entire pattern).
    \item Boresight gain (gain at $\mathrm{az}=0^\circ$, $\mathrm{el}=0^\circ$).
    \item Azimuth 3\,dB beamwidth (measured from the $\mathrm{el}=0^\circ$ cut).
    \item Elevation 3\,dB beamwidth (measured from the $\mathrm{az}=0^\circ$ cut).
    \item Minimum gain across the entire pattern.
    \item Average gain (power-averaged).
    \item Front-to-back ratio.
    \item First sidelobe levels (azimuth and elevation).
    \item Azimuth and elevation symmetry (RMS deviation).
    \item Angular coverage above threshold.
\end{itemize}

The result box displays the full analysis report.

\subsubsection{Card: Compare Multiple Patterns}

\begin{longtable}{L{3cm} L{3cm} L{7.2cm}}
\toprule
\textbf{Field} & \textbf{Widget Type} & \textbf{Description} \\
\midrule
\endhead
CSV Files & FilePicker (multi) & Multi-select CSV files for comparison. \\
Or Pattern Dir & FilePicker (directory) & Directory fallback — all CSVs in the directory are compared. \\
\bottomrule
\end{longtable}

\textbf{Button:} ``Compare'' (primary) — Runs \texttt{analyze\_single()} on each file, then \texttt{compare\_patterns()} to produce a side-by-side comparison table.

The result box displays the comparison table.

\subsubsection{Card: Pattern Set Comparison}

This card compares two sets of pattern files frequency-by-frequency, producing a percentage-difference table for each metric. It is modelled after the \texttt{scripts/append\_diff\_table.py} utility but integrated directly into the GUI.

\begin{longtable}{L{3.5cm} L{2.5cm} L{7.2cm}}
\toprule
\textbf{Field} & \textbf{Widget Type} & \textbf{Description} \\
\midrule
\endhead
Pattern Set A (files) & FilePicker (multi) & Pattern files for the first set (e.g.\ measured data). \\
Or Set A Dir & FilePicker (dir) & Directory fallback for Set A. \\
Pattern Set B (files) & FilePicker (multi) & Pattern files for the second set (e.g.\ fitted/generated data). \\
Or Set B Dir & FilePicker (dir) & Directory fallback for Set B. \\
\bottomrule
\end{longtable}

Both sets must contain the \textbf{same number of files}. Files are matched in sorted order, so ensure both sets use consistent naming.

\textbf{Button:} ``Compare Sets'' (primary) --- Analyses each file pair and computes percentage differences:
\[ \Delta\% = \frac{B - A}{|A|} \times 100 \]
for each of: peak gain, boresight gain, azimuth beamwidth, elevation beamwidth, and minimum gain.

The result box shows a formatted table with per-file differences plus summary statistics (average and maximum absolute differences across all frequency slices).

\subsubsection{Card: Sidelobe Mask Compliance}

This card tests a pattern against industry-standard sidelobe envelope masks. The card includes a documentation block explaining the three supported standards.

\paragraph{Supported Standards}
\begin{description}[style=nextline, leftmargin=2em]
    \item[\textbf{ITU-R S.580-6}] International Telecommunication Union recommendation for earth station antennas. The mask envelope is derived from the peak gain. Formula: $G(\theta) = 32 - 25\log_{10}(\theta)$ for the principal region (varies by angular range).
    \item[\textbf{ITU-R S.465-6}] ITU recommendation for small earth station antennas. The mask is parameterised by $D/\lambda$ (diameter-to-wavelength ratio).
    \item[\textbf{MIL-STD}] Military standard sidelobe envelope. Parameterised by peak gain and first sidelobe level.
\end{description}

\begin{longtable}{L{3cm} L{2.5cm} L{2cm} L{5.7cm}}
\toprule
\textbf{Field} & \textbf{Widget Type} & \textbf{Default} & \textbf{Description} \\
\midrule
\endhead
Pattern CSV & FilePicker & --- & The pattern file to test. \\
Mask Standard & Dropdown & ITU-R S.580-6 & The sidelobe mask standard. \\
Test Plane & Dropdown & az & The principal plane to test: \textbf{az} (azimuth at $\mathrm{el}=0^\circ$) or \textbf{el} (elevation at $\mathrm{az}=0^\circ$). \\
Peak Gain (dBi) & Text Entry & 35.0 & The peak gain value used in the mask formula. \\
\bottomrule
\end{longtable}

\textbf{Button:} ``Test Compliance'' (primary) — Evaluates the selected cut against the mask and reports:
\begin{itemize}
    \item Overall PASS/FAIL status.
    \item Number and location of violations.
    \item Minimum margin (dB below the mask).
    \item Worst-case violation angle and magnitude.
\end{itemize}

\subsubsection{Card: Parameter Fit}

This card enables reverse parameter fitting: given pre-existing antenna pattern files (CSV, DAT, or ODESSA \texttt{.tbl}), it derives the AntennaForge configuration parameters that best reproduce them. The fitted config can be saved as JSON or sent directly to the Generate page.

The fitter uses a multi-strategy approach: each run tries a different combination of physics features (sidelobes, ground reflection, etc.) and perturbed initial parameters. The optimizer employs a coarse-to-fine grid strategy (15\textdegree{} stride for global search, 3\textdegree{} for refinement), differential evolution for high-dimensional types (Dish, Array), cross-correlation cost functions with adaptive noise-floor weighting, and principal-plane-first fitting. A warm-start cache (\texttt{\~{}/.antennaforge/fit\_cache.json}) stores successful fits to seed future runs.

\begin{longtable}{L{3cm} L{2.5cm} L{2cm} L{5.7cm}}
\toprule
\textbf{Field} & \textbf{Widget Type} & \textbf{Default} & \textbf{Description} \\
\midrule
\endhead
Antenna Type & Dropdown & LPDA & Antenna model to fit against. Changing this auto-updates the recommended runs and target RMS to per-type defaults. \\
F Min (MHz) & Text Entry & 30.0 & Lower band edge of the reference antenna. \\
F Max (MHz) & Text Entry & 300.0 & Upper band edge of the reference antenna. \\
Pattern Files & FilePicker (multi) & --- & CSV, DAT, or ODESSA .tbl files. \\
Or Pattern Dir & FilePicker (dir) & --- & Directory of pattern files. \\
Optimize (SciPy) & Switch (toggle) & ON & When enabled and SciPy is installed, refines parameters via numerical optimisation. \\
Runs & Text Entry & (per-type) & Number of fitting strategy runs. Set to \textbf{0} for unlimited runs (use Stop button to halt). Recommended values: Monopole/Omni 2, LPDA/Horn/Panel 4, Dish 8, Array 10. \\
Target RMS & Text Entry & (per-type) & Target RMS error in dB to stop optimisation early. Per-type defaults: Monopole 0.5, Omni 1.0, LPDA/Horn/Panel 1.5, Dish 2.0, Array 2.5. \\
\bottomrule
\end{longtable}

\textbf{Buttons:}
\begin{itemize}
    \item ``Fit Parameters'' (primary) --- Runs the fitting algorithm in a background thread. Disabled while a fit is running.
    \item ``Stop'' (secondary) --- Appears while a fit is running. Signals the fitter to stop after the current run completes and return the best result found so far.
    \item ``Save Config JSON'' (secondary) --- Saves the fitted configuration to a JSON file. Enabled after a successful fit.
    \item ``Send to Generate'' (primary) --- Loads the fitted config onto the Generate page and navigates to it. Enabled after a successful fit.
\end{itemize}

The result box displays the fitting report with per-file RMS errors, the derived configuration, and which strategy produced the best result.

\subsubsection{Card: Batch Fitting}

Fit many pattern sets at once from a parent directory.

\begin{longtable}{L{3cm} L{2.5cm} L{2cm} L{5.7cm}}
\toprule
\textbf{Field} & \textbf{Widget Type} & \textbf{Default} & \textbf{Description} \\
\midrule
\endhead
Parent Dir & FilePicker (directory) & --- & Folder whose subfolders each contain one pattern set. \\
Output Dir & FilePicker (directory) & --- & Where fitted configs, generated patterns, and reports are saved. \\
Runs per set & Text Entry & 0 & Number of fitting runs per pattern set. \textbf{0} = auto (2\texttimes{} the recommended runs for the detected antenna type). \\
Target RMS & Text Entry & 0 & Target RMS error in dB. \textbf{0} = use per-type default. \\
\bottomrule
\end{longtable}

\textbf{Buttons:}
\begin{itemize}
    \item ``Run Batch Fit'' (primary) --- Launches the batch fitting pipeline in a background thread.
    \item ``Stop'' (secondary) --- Signals the batch pipeline to stop after the current set completes.
\end{itemize}

\textbf{Pipeline:}
\begin{enumerate}
    \item Discover subfolders under \emph{Parent Dir}.
    \item For each subfolder, auto-detect the antenna type and frequency range from \texttt{generation\_config.json} (if present) or by parsing filenames.
    \item Run the parameter fitter with the configured \emph{Runs per set} and \emph{Target RMS}.
    \item Save the resulting \texttt{fitted\_config.json} into the subfolder under \emph{Output Dir}.
    \item Generate patterns from the fitted configuration.
    \item Compare the original patterns against the generated patterns (percentage difference analysis).
    \item Produce a detailed per-set report section and append it to the grand summary.
\end{enumerate}

\textbf{Output:} A file \texttt{batch\_fit\_report.txt} is written to \emph{Output Dir} containing a grand summary table (one row per pattern set with type, frequency range, best RMS, and pass/fail status) followed by detailed per-set sections with individual file RMS values and fitted parameter summaries.

% ── OPERATIONS PAGE ──────────────────────────────────────────────
\subsection{Operations Page}
\label{sec:page_operations}

The Operations page provides pattern arithmetic, three-axis rotation, and interpolation/resampling.

\subsubsection{Card: Pattern Arithmetic}

\begin{longtable}{L{3.5cm} L{2.5cm} L{2cm} L{5.2cm}}
\toprule
\textbf{Field} & \textbf{Widget Type} & \textbf{Default} & \textbf{Description} \\
\midrule
\endhead
Operation & Dropdown & Add (power) & The arithmetic operation to perform. Options described below. \\
Pattern A & FilePicker & --- & First input CSV. \\
Pattern B & FilePicker & --- & Second input CSV (not used for Scale, Max Envelope, or Average). \\
Multi Files & FilePicker (multi) & --- & Multi-select for Max Envelope and Average operations. \\
Scale Amount (dB) & Text Entry & 0.0 & Constant dB offset for the Scale operation. \\
Output CSV & FilePicker (save) & --- & Output file path. \\
\bottomrule
\end{longtable}

\paragraph{Arithmetic Operations}
\begin{description}[style=nextline, leftmargin=2em]
    \item[\textbf{Add (power)}] Converts both patterns to linear power, sums element-wise, converts back to dB. Models coherent power combination.
    \item[\textbf{Subtract (dB)}] Computes the dB difference: $A_{\text{dB}} - B_{\text{dB}}$ at each (az, el) point. Useful for isolation or interference analysis.
    \item[\textbf{Multiply (dB+dB)}] Adds the two patterns in dB (multiplication in linear). Implements the pattern multiplication theorem for array analysis.
    \item[\textbf{Scale (dB offset)}] Adds a constant offset to all gain values: $G_{\text{out}} = G_{\text{in}} + \Delta_{\text{dB}}$. Useful for modelling cable loss or amplifier gain.
    \item[\textbf{Max Envelope}] Computes the point-by-point maximum across multiple patterns. Useful for worst-case interference analysis.
    \item[\textbf{Average (power)}] Converts to linear, averages across all files, converts back to dB. Useful for statistical pattern analysis.
\end{description}

\textbf{Button:} ``Run Arithmetic'' (primary).

\subsubsection{Card: Pattern Rotation}

\begin{longtable}{L{3.5cm} L{2.5cm} L{1.5cm} L{2cm} L{3.7cm}}
\toprule
\textbf{Field} & \textbf{Widget Type} & \textbf{Default} & \textbf{Range} & \textbf{Description} \\
\midrule
\endhead
Input CSV & FilePicker & --- & --- & Single CSV to rotate. \\
Or Input Dir & FilePicker (dir) & --- & --- & Directory of CSVs. \\
Yaw / Az Rotate ($^\circ$) & Slider & 0 & $-180$ to $180$ & Azimuth rotation. \\
Pitch / El Tilt ($^\circ$) & Slider & 0 & $-90$ to $90$ & Elevation tilt. \\
Roll ($^\circ$) & Slider & 0 & $-180$ to $180$ & Roll rotation. \\
Output CSV / Dir & FilePicker (save) & --- & --- & Output path. \\
\bottomrule
\end{longtable}

The rotation algorithm:
\begin{enumerate}
    \item Constructs a $3 \times 3$ rotation matrix from the three Euler angles (yaw, pitch, roll).
    \item Converts each output (az, el) grid point to Cartesian coordinates.
    \item Applies the inverse rotation matrix to find the source direction.
    \item Converts back to (az, el) and interpolates the gain from the original pattern using bilinear interpolation (or bicubic via SciPy's \texttt{RegularGridInterpolator} when available).
\end{enumerate}

\textbf{Button:} ``Rotate'' (primary).

\subsubsection{Card: Interpolation \& Resampling}

\begin{longtable}{L{3.5cm} L{2.5cm} L{2cm} L{5.2cm}}
\toprule
\textbf{Field} & \textbf{Widget Type} & \textbf{Default} & \textbf{Description} \\
\midrule
\endhead
Mode & Dropdown & Freq Interpolate (2 files) & Interpolation mode. Options described below. \\
Low-Freq CSV & FilePicker & --- & Lower frequency pattern (2-file mode). \\
High-Freq CSV & FilePicker & --- & Higher frequency pattern (2-file mode). \\
Target Freq (MHz) & Text Entry & (blank) & Target frequency or comma-separated list. \\
Multi Files & FilePicker (multi) & --- & Multi-select for multi-freq interpolation. \\
New Az Step ($^\circ$) & Text Entry & 1.0 & New azimuth resolution (resample mode). \\
New El Step ($^\circ$) & Text Entry & 1.0 & New elevation resolution (resample mode). \\
\bottomrule
\end{longtable}

\paragraph{Interpolation Modes}
\begin{description}[style=nextline, leftmargin=2em]
    \item[\textbf{Freq Interpolate (2 files)}] Linear interpolation between two patterns at different frequencies to estimate the pattern at an intermediate frequency. The frequency is extracted from each filename.
    \item[\textbf{Multi-Freq Interpolate}] Given multiple frequency-stamped pattern files, interpolates to one or more target frequencies using piecewise linear interpolation between the nearest bracketing files.
    \item[\textbf{Resample Grid}] Resamples a pattern to a new angular grid resolution using SciPy's \texttt{RectBivariateSpline} (bicubic) when available, or bilinear interpolation as fallback.
\end{description}

\textbf{Button:} ``Interpolate'' (primary).

% ── RF CALCS PAGE ────────────────────────────────────────────────
\subsection{RF Calcs Page}
\label{sec:page_rfcalc}

The RF Calcs page provides EIRP calculation, power density computation, free-space path loss, and link budget analysis.

\subsubsection{Card: EIRP Calculator}

\begin{longtable}{L{3cm} L{2.5cm} L{1.5cm} L{6.2cm}}
\toprule
\textbf{Field} & \textbf{Widget Type} & \textbf{Default} & \textbf{Description} \\
\midrule
\endhead
Pattern CSV & FilePicker & --- & The antenna gain pattern. \\
Tx Power (dBm) & Text Entry & 30.0 & Transmitter output power (30\,dBm = 1\,W). \\
Cable Loss (dB) & Text Entry & 2.0 & Feedline/connector attenuation. \\
Distance (m) & Text Entry & 1000.0 & Distance for power density calculation. \\
\bottomrule
\end{longtable}

\paragraph{EIRP Formula}
\begin{equation}
\mathrm{EIRP}(\mathrm{az}, \mathrm{el}) = P_{\text{tx,dBm}} - L_{\text{cable,dB}} + G(\mathrm{az}, \mathrm{el})_{\text{dBi}}
\end{equation}

\paragraph{Power Density Formula}
\begin{equation}
\mathrm{PD}(\mathrm{az}, \mathrm{el}) = \frac{\mathrm{EIRP}_{\text{linear}}(\mathrm{az}, \mathrm{el})}{4\pi d^2} \quad [\text{W/m}^2]
\end{equation}

\textbf{Buttons:} ``Calculate EIRP'' (primary), ``Calculate Power Density'' (secondary).

\subsubsection{Card: Free-Space Path Loss}

\begin{longtable}{L{3cm} L{2.5cm} L{1.5cm} L{6.2cm}}
\toprule
\textbf{Field} & \textbf{Widget Type} & \textbf{Default} & \textbf{Description} \\
\midrule
\endhead
Distance (km) & Text Entry & 100.0 & Link distance in kilometres. \\
Frequency (MHz) & Text Entry & 1000.0 & Operating frequency in MHz. \\
\bottomrule
\end{longtable}

\paragraph{FSPL Formula}
\begin{equation}
\mathrm{FSPL}_{\text{dB}} = 20\log_{10}(d_{\text{km}}) + 20\log_{10}(f_{\text{MHz}}) + 32.44
\end{equation}

\textbf{Button:} ``Calculate FSPL'' (primary).

\subsubsection{Card: Link Budget Analysis}

\begin{longtable}{L{3.5cm} L{2.5cm} L{1.5cm} L{5.7cm}}
\toprule
\textbf{Field} & \textbf{Widget Type} & \textbf{Default} & \textbf{Description} \\
\midrule
\endhead
Mode & Dropdown & Boresight Link Margin & Calculation mode: \textbf{Boresight Link Margin} or \textbf{Full Angular Map}. \\
Tx Pattern CSV & FilePicker & --- & Transmit antenna pattern. \\
Rx Pattern CSV & FilePicker & --- & Receive antenna pattern. \\
Distance (km) & Text Entry & 100.0 & Link distance. \\
Frequency (MHz) & Text Entry & 1000.0 & Operating frequency. \\
Tx Power (dBm) & Text Entry & 30.0 & Transmitter power. \\
Tx Cable Loss (dB) & Text Entry & 2.0 & Tx-side cable/connector loss. \\
Rx Cable Loss (dB) & Text Entry & 1.0 & Rx-side cable/connector loss. \\
Rx Sensitivity (dBm) & Text Entry & $-90.0$ & Receiver sensitivity threshold. Only used in Boresight Link Margin mode. \\
\bottomrule
\end{longtable}

\paragraph{Link Budget Formula}
\begin{equation}
P_{\text{rx}} = P_{\text{tx}} - L_{\text{tx}} + G_{\text{tx}}(\mathrm{az}_{\text{tx}}) - \mathrm{FSPL} + G_{\text{rx}}(\mathrm{az}_{\text{rx}}) - L_{\text{rx}}
\end{equation}

\paragraph{Link Margin}
\begin{equation}
\text{Margin} = P_{\text{rx,boresight}} - P_{\text{sensitivity}}
\end{equation}

\paragraph{Modes}
\begin{description}[style=nextline, leftmargin=2em]
    \item[\textbf{Boresight Link Margin}] Computes received power using boresight (peak) gains of both antennas and reports the margin above receiver sensitivity.
    \item[\textbf{Full Angular Map}] Computes received power for every combination of Tx and Rx orientations. Outputs a 2D CSV of received power indexed by transmit azimuth and receive azimuth.
\end{description}

\textbf{Button:} ``Calculate Link Budget'' (primary).

% ── REPORTS PAGE ─────────────────────────────────────────────────
\subsection{Reports Page}
\label{sec:page_reports}

The Reports page generates comprehensive PDF or \LaTeX\ reports from a directory of pattern CSVs.

\subsubsection{Card: Input}

\begin{longtable}{L{3cm} L{2.5cm} L{1.5cm} L{6.2cm}}
\toprule
\textbf{Field} & \textbf{Widget Type} & \textbf{Default} & \textbf{Description} \\
\midrule
\endhead
Pattern Directory & FilePicker (dir) & --- & Source directory containing pattern CSVs. \\
Output Format & Dropdown & PDF & Report format: \textbf{PDF} or \textbf{LaTeX (.tex)}. \\
Report Title & Text Entry & (blank) & Custom report title. If blank, an automatic title is generated from the antenna type and frequency band. \\
\bottomrule
\end{longtable}

\subsubsection{Card: Presets}

\begin{longtable}{L{2.5cm} L{10.7cm}}
\toprule
\textbf{Preset} & \textbf{Sections Enabled} \\
\midrule
\endhead
\texttt{default} & Cover, TOC, Summary Table, Gain vs Freq, BW vs Freq, Overlay Cuts, Per-freq Heatmap, Polar, Stats, Detailed, Sidelobe. \\
\texttt{full} & All 13 sections enabled (including 3D Surface and Cross-Pol). \\
\texttt{minimal} & Cover, TOC, Summary Table, plus per-freq Heatmap and per-freq Statistics. \\
\texttt{quick} & Global pages only (Cover, TOC, Summary Table, trend plots, overlay cuts). No per-frequency detail. \\
\bottomrule
\end{longtable}

\subsubsection{Card: Report Sections}
Thirteen toggle switches arranged in two columns:

\begin{longtable}{L{3.5cm} L{2cm} L{1.5cm} L{6.2cm}}
\toprule
\textbf{Section} & \textbf{Group} & \textbf{Default} & \textbf{Description} \\
\midrule
\endhead
Cover Page & global & ON & Title page with project metadata, author, classification. \\
Table of Contents & global & ON & Auto-generated TOC. \\
Summary Table & global & ON & Frequency, peak gain, beamwidths, and front-to-back ratio for all slices. \\
Gain vs Freq & global & ON & Line plot of peak gain across frequency. \\
BW vs Freq & global & ON & Line plot of azimuth and elevation beamwidth across frequency. \\
Overlay Cuts & global & ON & Overlaid azimuth cuts from all frequency slices. \\
Per-freq Heatmap & per\_freq & ON & Full 2D heatmap for each frequency slice. \\
Per-freq Polar & per\_freq & ON & Polar plot for each frequency slice. \\
Per-freq 3D Surface & per\_freq & OFF & 3D surface for each frequency (computationally expensive). \\
Per-freq Statistics & per\_freq & ON & Detailed statistics table for each slice. \\
Per-freq Detailed & per\_freq & ON & Comprehensive per-slice analysis (beamwidth, sidelobes, symmetry). \\
Per-freq Sidelobe & per\_freq & ON & Sidelobe analysis for each frequency slice. \\
Include Cross-Pol & option & OFF & Include cross-polarisation patterns in the report (if \texttt{\_xpol.csv} files exist). \\
\bottomrule
\end{longtable}

\subsubsection{Card: Report Options}

\begin{longtable}{L{3.5cm} L{2cm} L{1.5cm} L{6.2cm}}
\toprule
\textbf{Field} & \textbf{Widget Type} & \textbf{Default} & \textbf{Description} \\
\midrule
\endhead
Dynamic Range (dB) & Text Entry & 60 & The dB range of the heatmap colour scale. Gain values below peak $-$ dynamic range are plotted at the minimum colour. \\
Polar Range (dB) & Text Entry & 40 & The dB range displayed in polar plots. \\
\bottomrule
\end{longtable}

\subsubsection{Card: Cover Page Fields}

\begin{longtable}{L{2.5cm} L{2cm} L{3.5cm} L{5.2cm}}
\toprule
\textbf{Field} & \textbf{Widget Type} & \textbf{Default} & \textbf{Description} \\
\midrule
\endhead
Subtitle & Text Entry & Antenna Pattern Report & Cover page subtitle. \\
Author & Text Entry & (blank) & Author name. \\
Organisation & Text Entry & (blank) & Organisation or company name. \\
Project & Text Entry & (blank) & Project identifier. \\
Document ID & Text Entry & (blank) & Document tracking number. \\
Revision & Text Entry & (blank) & Document revision string. \\
Date & Text Entry & (blank) & Cover date. Blank = today's date. \\
\bottomrule
\end{longtable}

\subsubsection{Card: Classification Markings}

\begin{longtable}{L{3.5cm} L{2.5cm} L{2cm} L{5.2cm}}
\toprule
\textbf{Field} & \textbf{Widget Type} & \textbf{Default} & \textbf{Description} \\
\midrule
\endhead
Classification Level & Dropdown & UNCLASSIFIED & Security classification applied to every page header and footer. Options: UNCLASSIFIED, CUI, CONFIDENTIAL, SECRET, TOP~SECRET, TOP~SECRET//SCI. \\
Caveats / Releasability & Text Entry & (blank) & Additional handling caveats (e.g., \texttt{//NOFORN}, \texttt{//REL TO USA, GBR}). \\
\bottomrule
\end{longtable}

\textbf{Button:} ``Generate Report'' (success/green).

% ── HOW TO PAGE ──────────────────────────────────────────────────
\subsection{How To Page}
\label{sec:page_howto}

The How To page is a built-in reference guide consisting of nine read-only information cards. It contains no interactive widgets. The topics covered are:

\begin{enumerate}
    \item \textbf{Quick Start} — A six-step workflow from generation to reporting.
    \item \textbf{Generate \textemdash\ Creating Patterns} \textemdash\ Core parameters, beamwidth, frequency, polarisation, physics features, sidelobes, output, and batch processing.
    \item \textbf{Visualize \textemdash\ Plotting Patterns} \textemdash\ Input sources, plot types, caching, and navigation.
    \item \textbf{Analysis \textemdash\ Inspecting Patterns} \textemdash\ Single analysis, comparison, and mask compliance with ITU/MIL references.
    \item \textbf{Operations \textemdash\ Pattern Math} \textemdash\ Arithmetic (six operations), rotation (yaw/pitch/roll), and interpolation (three modes).
    \item \textbf{RF Calcs \textemdash\ Link Budget \& EIRP} \textemdash\ EIRP, power density, FSPL, and link budget.
    \item \textbf{Reports \textemdash\ PDF \& LaTeX} \textemdash\ Presets, section toggles, and cover page.
    \item \textbf{Settings \textemdash\ Appearance \& Config} \textemdash\ Theme, fonts, colours, and config I/O.
    \item \textbf{Tips \& Tricks} — Tooltip hints, slider/entry sync, caching, and shortcuts.
\end{enumerate}

% ── SETTINGS PAGE ────────────────────────────────────────────────
\subsection{Settings Page}
\label{sec:page_settings}

The Settings page provides theme customisation, configuration file management, and defaults review.

\subsubsection{Card: Appearance}

\begin{longtable}{L{3cm} L{2.5cm} L{2cm} L{5.7cm}}
\toprule
\textbf{Field} & \textbf{Widget Type} & \textbf{Default} & \textbf{Description} \\
\midrule
\endhead
Colour Theme & Dropdown & dark & Global appearance mode: \textbf{dark}, \textbf{light}, or \textbf{system} (follows OS setting). \\
UI Font Family & Dropdown & Segoe UI & Main UI typeface. Curated list of approximately 100+ system sans-serif, decorative, and display fonts. \\
Mono Font Family & Dropdown & Consolas & Monospace typeface for code and result boxes. Curated list of approximately 22 monospace fonts. \\
Base Font Size & Dropdown & 12 & Base font size in points. Options: 9, 10, 11, 12, 13, 14, 15, 16. All other font sizes are derived from this base. \\
\bottomrule
\end{longtable}

\textbf{Buttons:}
\begin{itemize}
    \item ``Apply Fonts'' (primary) — Persists the font settings and rebuilds the entire UI.
    \item ``Reset Fonts'' (secondary) — Restores font settings to defaults and rebuilds.
\end{itemize}

\subsubsection{Card: Colour Customization}
\label{sec:colour_custom}

All 27 colours in the theme palette can be individually customised. Each colour is displayed as a row with a label, light-mode and dark-mode hex entries, and a live colour swatch that updates on each keystroke. The colours are organised in six groups matching the theme system (Section~\ref{sec:theme}): Sidebar (5), Main (5), Text (3), Accent (3), Status (4), Inputs (3).

\textbf{Buttons:}
\begin{itemize}
    \item ``Apply Colours'' (primary) — Applies colour overrides, persists to preferences file, and rebuilds the UI.
    \item ``Reset to Defaults'' (secondary) — Restores the factory colour palette.
    \item ``Export Theme JSON'' (secondary) — Opens a save-as dialog to export the current colour theme as a JSON file.
    \item ``Import Theme JSON'' (secondary) — Opens a file dialog to import a colour theme JSON and populate the entry fields.
\end{itemize}

\subsubsection{Card: Configuration Files}

\begin{longtable}{L{3cm} L{3cm} L{7.2cm}}
\toprule
\textbf{Field} & \textbf{Widget Type} & \textbf{Description} \\
\midrule
\endhead
Load Config JSON & FilePicker & Select a previously saved configuration JSON file. \\
\bottomrule
\end{longtable}

\textbf{Buttons:}
\begin{itemize}
    \item ``Load Config'' (primary) — Loads the selected JSON file and displays its contents in the result box.
    \item ``Save Current Config'' (secondary) — Saves the current Generate page configuration to a JSON file via a save-as dialog.
\end{itemize}

\subsubsection{Card: Current Default Configuration}

\textbf{Button:} ``Show Defaults'' (secondary) — Displays the complete \texttt{DEFAULT\_CONFIG} dictionary as formatted JSON in the result box.

% ===================================================================
%  PART VII — CLI REFERENCE
% ===================================================================
\newpage
\section{Command-Line Interface}
\label{sec:cli}

\subsection{Interactive CLI Menu}
The interactive CLI mode provides a state-machine-driven menu with GoBack support. Launch with:
\begin{lstlisting}[language=bash]
python -m antenna_tool
\end{lstlisting}

\subsubsection{Main Menu Options}

\begin{longtable}{cl L{8cm}}
\toprule
\textbf{Key} & \textbf{Option} & \textbf{Description} \\
\midrule
\endhead
1 & Generate & Launches the 8-step generation wizard. \\
2 & Graph & Plot patterns (heatmap, cuts, polar, 3D, gain/BW vs freq, overlay, animation). \\
3 & Report & Generate PDF or \LaTeX\ reports. \\
4 & Analyze & Single analysis, comparison, or batch analysis. \\
5 & Arithmetic & Pattern arithmetic operations. \\
6 & Rotate & Three-axis pattern rotation. \\
7 & Mask Test & Sidelobe mask compliance testing. \\
8 & EIRP & EIRP and power density calculation. \\
9 & Link Budget & Link budget analysis. \\
A & Interpolate & Frequency interpolation and angular resampling. \\
B & Batch & Batch processing from recipe files. \\
Q & Quit & Exit the application. \\
\bottomrule
\end{longtable}

\subsubsection{Generation Wizard (8 Steps)}
The generation wizard uses a state machine with breadcrumb navigation:

\begin{enumerate}
    \item \textbf{Core} — Antenna type, gain, F/B ratio, antenna-specific prompts.
    \item \textbf{Beamwidth} — Azimuth and elevation beamwidths.
    \item \textbf{Frequency} — Band limits, number of slices, spacing, reference frequency.
    \item \textbf{Angular} — Azimuth and elevation step sizes.
    \item \textbf{Polarization} — Polarisation mode and custom angle.
    \item \textbf{Features} — Physics feature toggles and configuration.
    \item \textbf{Output} — Output directory selection.
    \item \textbf{Summary} — Review all settings, confirm or go back.
\end{enumerate}

Each step supports ``b'' or ``back'' to return to the previous step.

\subsection{Direct CLI Arguments}
For scriptable operation, pass arguments directly:
\begin{lstlisting}[language=bash,caption={CLI argument examples}]
# Basic generation
python -m antenna_tool --type LPDA --max-gain 10 --az-bw 65 --el-bw 70 \
    --f-min 30 --f-max 300 --slices 10 --output ./patterns

# With features
python -m antenna_tool --type Panel --max-gain 15 --az-bw 65 --el-bw 30 \
    --f-min 800 --f-max 2500 --slices 20 --ftb 25 \
    --xpol --ground --no-breakup

# Batch mode
python -m antenna_tool --batch recipe.json

# Batch with image generation and custom output folder
python -m antenna_tool --batch recipe.json --batch-images \
    --batch-output ./batch_results

# Batch with images and PDF reports (minimal preset)
python -m antenna_tool --batch recipe.json --batch-images \
    --batch-report --batch-report-preset minimal \
    --batch-output ./batch_results

# Report generation
python -m antenna_tool --report ./patterns

# LaTeX report
python -m antenna_tool --latex-report ./patterns
\end{lstlisting}

% ===================================================================
%  PART VIII — CORE ENGINE DETAILS
% ===================================================================
\newpage
\section{Core Engine Architecture}
\label{sec:engine}

\subsection{Pattern Computation Pipeline}
The \texttt{core/engine.py} module drives pattern generation. The computation pipeline for each frequency slice is:

\begin{enumerate}
    \item \textbf{Frequency-dependent beamwidth:} Compute the effective azimuth and elevation beamwidths at the current frequency using the selected decay mode (exponent, percentage, or three-point).
    \item \textbf{VSWR rolloff:} Compute the gain reduction at this frequency based on its position within the band. The VSWR rolloff feature applies a steep gain drop at the last frequency slice to model realistic impedance mismatch effects. This is visible in generated patterns and trend plots.
    \item \textbf{Asymmetry (optional):} Compute azimuth squint and elevation tilt offsets.
    \item \textbf{Grid computation:} For each (azimuth, elevation) point:
    \begin{enumerate}[label=\alph*.]
        \item Apply asymmetry offsets.
        \item Compute the antenna model's point gain (linear).
        \item Apply the sigmoid front-fade function.
        \item Add the sidelobe envelope (radial or per-plane).
        \item Apply pattern breakup ripple (if enabled).
        \item Apply ground reflection interference (if enabled).
        \item Apply polarisation weighting.
        \item Add noise.
    \end{enumerate}
    \item \textbf{Cross-polarisation (optional):} Compute the cross-pol gain at each point.
    \item \textbf{Output:} Convert to dBi and write to CSV.
\end{enumerate}

\subsection{Vectorised vs Scalar Paths}
When NumPy is available, the engine uses fully vectorised 2D array operations (\texttt{\_compute\_pattern\_numpy}), processing the entire azimuth$\times$elevation grid at once. This is 10--100$\times$ faster than the pure-Python scalar fallback. SciPy enables additional optimisations such as Bessel function evaluation for the Dish model.

Each antenna type has a dedicated vectorised gain function. The Array type uses \texttt{\_vectorized\_array\_gain()}, which computes element patterns, array factors (linear, planar, or circular), mutual coupling, and element errors entirely via NumPy broadcasting --- weights, coupling corrections, and error perturbations are computed once per call and the array factor is evaluated across the full angle grid in a single vectorised pass. Without this path, Array patterns fall through to the per-point scalar loop, which is prohibitively slow for large element counts (e.g., 16-element arrays with mutual coupling).

\subsection{Pattern Math Library}
The \texttt{core/pattern\_math.py} module provides all mathematical building blocks:

\begin{longtable}{L{4.5cm} L{8.7cm}}
\toprule
\textbf{Function} & \textbf{Purpose} \\
\midrule
\endhead
\texttt{front\_fade\_scalar} & Smooth sigmoid front-hemisphere weighting. \\
\texttt{bw\_to\_exponent} & Converts 3\,dB beamwidth to $\cos^n$ exponent. \\
\texttt{compute\_freq\_dependent\_bw} & Beamwidth at a given frequency (3 modes). \\
\texttt{compute\_sidelobe\_envelope} & Sidelobe gain at a given off-axis angle. \\
\texttt{compute\_sidelobe\_radial} & 2D radial sidelobe envelope. \\
\texttt{apply\_pattern\_breakup} & Deterministic ripple at high angles. \\
\texttt{apply\_ground\_reflection} & 2-ray ground interference model. \\
\texttt{compute\_cross\_pol} & Cross-pol gain computation. \\
\texttt{apply\_vswr\_rolloff} & Band-edge gain reduction. \\
\texttt{apply\_asymmetry} & Pointing offset and perturbation. \\
\bottomrule
\end{longtable}

\subsection{Parameter Fitting Module}
The \texttt{core/fitting.py} module provides reverse parameter fitting --- deriving AntennaForge configuration parameters from reference (measured or legacy) patterns.

\begin{longtable}{L{4.5cm} L{8.7cm}}
\toprule
\textbf{Function} & \textbf{Purpose} \\
\midrule
\endhead
\texttt{fit\_parameters} & Analyses reference pattern files to extract gain, beamwidth, F/B ratio, and sidelobe levels; seeds an AntennaForge config; optionally refines via SciPy optimisation. \\
\texttt{fit\_parameters\_multirun} & Runs multiple fitting strategies with perturbed initial conditions and returns the best result. Supports unlimited runs (\texttt{n\_runs=0}) with a \texttt{stop\_event} for cancellation. \\
\texttt{recommended\_runs} & Returns the recommended number of runs and a reason string for a given antenna type. \\
\texttt{type\_target\_rms} & Returns the per-type auto-stop RMS threshold (dB). \\
\texttt{batch\_fit} & Fits many pattern sets from subfolders of a parent directory. For each set: auto-detects type and frequency range (from \texttt{generation\_config.json} or filename parsing), runs \texttt{fit\_parameters\_multirun}, saves \texttt{fitted\_config.json}, generates patterns from the fitted config, compares against originals, and produces a detailed report. Supports progress callbacks and a stop event for cancellation. \\
\bottomrule
\end{longtable}

\paragraph{Optimization Pipeline}
The fitting process uses a multi-stage pipeline:
\begin{enumerate}
    \item \textbf{Measurement Seed:} Extract gain, beamwidth, F/B ratio, and sidelobe levels from the reference patterns to build an initial configuration.
    \item \textbf{Strategy Generation:} For each run, generate a strategy with a different combination of physics features (sidelobes on/off, ground reflection on/off, etc.) and perturbed parameters.
    \item \textbf{Principal-Plane Fitting (Stage 0):} Fit azimuth and elevation cuts before the full 2-D pattern optimisation.
    \item \textbf{Coarse Optimisation (Stage 1):} Evaluate the cost function on a 15\textdegree{} grid stride for fast global search. Uses differential evolution for high-dimensional types (Dish, Array) or sensitivity-screened Nelder-Mead for simpler types.
    \item \textbf{Fine Optimisation (Stage 2):} Polish on a 3\textdegree{} grid stride via Nelder-Mead.
    \item \textbf{Ensemble Averaging:} Top-3 results are weighted-averaged (by $1/\text{score}$) and accepted if better than the single best.
    \item \textbf{Cache:} Successful fits are saved to \texttt{\~{}/.antennaforge/fit\_cache.json} (50 entries per type) to seed future runs.
\end{enumerate}

\paragraph{Cost Function}
The cost function blends normalised cross-correlation (60\%) with weighted percentage error (40\%), using per-type cost profiles that emphasise the beam core or far-out regions differently. Adaptive noise-floor weighting estimates the measurement noise from low-gain regions and adjusts the weighting of far-out angles accordingly.

The fitter supports CSV, DAT, and ODESSA \texttt{.tbl} input formats. ODESSA tables are read via \texttt{core/io.read\_odessa\_tbl()}, which parses multi-frequency slice data from a single file.

\subsection{I/O Utilities}
The \texttt{core/io.py} module provides additional I/O beyond basic CSV read/write:
\begin{itemize}
    \item \texttt{read\_odessa\_tbl(filepath)} --- Reads an ODESSA \texttt{.tbl} file containing multiple frequency slices. Returns a dictionary with keys: \texttt{az\_deg}, \texttt{el\_deg}, \texttt{freqs\_mhz}, and \texttt{slices} (a list of 2-D gain patterns).
    \item \texttt{make\_timestamped\_dir(base\_path)} --- Creates a timestamped output directory.
    \item \texttt{make\_graph\_output\_dir(source\_path)} --- Creates a timestamped graphs directory alongside the source data.
\end{itemize}

% ===================================================================
%  PART IX — ANALYSIS MODULE
% ===================================================================
\newpage
\section{Analysis Module}
\label{sec:analysis}

The analysis module (\texttt{analysis/analyzer.py}) provides comprehensive pattern characterisation.

\subsection{Single Pattern Analysis}
The \texttt{analyze\_single()} function computes:
\begin{itemize}
    \item \textbf{Peak gain:} Maximum gain value anywhere in the pattern.
    \item \textbf{Boresight gain:} Gain at $(\mathrm{az}=0^\circ, \mathrm{el}=0^\circ)$.
    \item \textbf{Azimuth 3\,dB beamwidth:} Measured from the elevation=$0^\circ$ cut by finding the $-3$\,dB points on either side of the peak.
    \item \textbf{Elevation 3\,dB beamwidth:} Measured from the azimuth=$0^\circ$ cut.
    \item \textbf{Minimum gain:} Lowest gain in the entire pattern.
    \item \textbf{Average gain:} Power-averaged gain across all angles.
    \item \textbf{Front-to-back ratio:} Difference between boresight gain and the gain at $180^\circ$ azimuth.
    \item \textbf{First sidelobe levels:} Detected in both azimuth and elevation cuts.
    \item \textbf{Symmetry:} RMS deviation between left/right (azimuth) and upper/lower (elevation) pattern halves.
    \item \textbf{Angular coverage:} Solid angle above a gain threshold.
\end{itemize}

\subsection{Pattern Comparison}
The \texttt{compare\_patterns()} function tabulates side-by-side statistics for multiple patterns, enabling rapid comparison across frequency slices or antenna configurations.

\subsection{In-Memory Analysis}
The \texttt{analyze\_from\_data(az\_angles, el\_angles, data\_2d)} function performs the same analysis as \texttt{analyze\_single()} but operates on in-memory data arrays rather than loading from a file. This is used by the parameter fitting module to evaluate generated patterns without intermediate disk I/O.

\subsection{Beamwidth Measurement Algorithm}
The \texttt{measure\_beamwidth()} function:
\begin{enumerate}
    \item Finds the peak gain in the 1D cut.
    \item Computes the $-3$\,dB threshold.
    \item Scans outward from the peak in both directions to find the crossing points.
    \item Uses linear interpolation between adjacent samples for sub-degree accuracy.
    \item Returns the angular width between the two crossing points.
\end{enumerate}

% ===================================================================
%  PART X — PATTERN OPERATIONS
% ===================================================================
\newpage
\section{Pattern Operations}
\label{sec:operations}

\subsection{Arithmetic Operations}
Six arithmetic operations are provided in \texttt{core/pattern\_ops.py}:

\begin{longtable}{L{3cm} L{10.2cm}}
\toprule
\textbf{Operation} & \textbf{Algorithm} \\
\midrule
\endhead
Add (power) & $G_{\text{out}} = 10\log_{10}\bigl(10^{G_A/10} + 10^{G_B/10}\bigr)$ \\
Subtract (dB) & $G_{\text{out}} = G_A - G_B$ (element-wise in dB) \\
Multiply (dB+dB) & $G_{\text{out}} = G_A + G_B$ (dB addition = linear multiplication) \\
Scale (offset) & $G_{\text{out}} = G_{\text{in}} + \Delta_{\text{dB}}$ \\
Max Envelope & $G_{\text{out}}(i,j) = \max\bigl(G_1(i,j), G_2(i,j), \ldots, G_N(i,j)\bigr)$ \\
Average (power) & $G_{\text{out}} = 10\log_{10}\!\left(\frac{1}{N}\sum_{k=1}^{N} 10^{G_k/10}\right)$ \\
\bottomrule
\end{longtable}

All operations validate that input patterns share the same angular grid before proceeding.

\subsection{Pattern Rotation}
The rotation system (\texttt{core/rotation.py}) applies three-axis Euler rotation:
\begin{enumerate}
    \item Constructs a $3\times3$ rotation matrix $\mathbf{R}$ from yaw ($\psi$), pitch ($\theta$), and roll ($\phi$).
    \item For each output grid point $(\mathrm{az}_o, \mathrm{el}_o)$:
    \begin{enumerate}[label=\alph*.]
        \item Convert to Cartesian: $(x, y, z) = (\cos\mathrm{el}\cos\mathrm{az},\;\cos\mathrm{el}\sin\mathrm{az},\;\sin\mathrm{el})$.
        \item Apply the inverse rotation: $(x', y', z') = \mathbf{R}^{-1}(x, y, z)$.
        \item Convert back to spherical: $(\mathrm{az}_s, \mathrm{el}_s) = (\arctan2(y', x'),\;\arcsin(z'))$.
        \item Interpolate the gain at $(\mathrm{az}_s, \mathrm{el}_s)$ from the original pattern.
    \end{enumerate}
\end{enumerate}

Interpolation uses bilinear interpolation in the pure-Python path, or SciPy's \texttt{RegularGridInterpolator} (bicubic) when available.

\subsection{Interpolation \& Resampling}
Three interpolation modes are provided in \texttt{core/interpolation.py}:

\begin{description}[style=nextline, leftmargin=2em]
    \item[\textbf{Frequency Interpolation (2 files)}] Linear interpolation in dB between two patterns at different frequencies: $G(f) = G_{\text{lo}} + (G_{\text{hi}} - G_{\text{lo}}) \cdot \frac{f - f_{\text{lo}}}{f_{\text{hi}} - f_{\text{lo}}}$.
    \item[\textbf{Multi-Frequency Interpolation}] Given $N$ patterns at known frequencies, interpolates to arbitrary target frequencies using piecewise linear interpolation between the nearest available files.
    \item[\textbf{Angular Resampling}] Resamples a pattern to a new azimuth/elevation grid resolution. Uses SciPy's \texttt{RectBivariateSpline} for bicubic interpolation when available.
\end{description}

% ===================================================================
%  PART XI — REPORTING
% ===================================================================
\newpage
\section{Reporting}
\label{sec:reporting}

AntennaForge generates three types of reports from a directory of pattern CSVs.

\subsection{PDF Reports}
The PDF report generator (\texttt{reporting/report.py}) uses Matplotlib's \texttt{PdfPages} backend to produce a multi-page document. The report structure follows the 13 configurable sections described in Section~\ref{sec:page_reports}.

Key features:
\begin{itemize}
    \item Cover page with project metadata and classification banner.
    \item Auto-generated table of contents.
    \item Summary table with peak gain, beamwidths, and F/B ratio for every frequency slice.
    \item Trend plots (gain vs frequency, beamwidth vs frequency).
    \item Overlay cuts for cross-frequency comparison.
    \item Per-frequency heatmaps, polar plots, 3D surfaces, statistics, and sidelobe analysis.
    \item Optional cross-polarisation sections.
    \item Classification markings on every page (header and footer).
\end{itemize}

\subsection{\LaTeX\ Reports}
The \LaTeX\ report generator (\texttt{reporting/latex\_report.py}) produces a standalone \texttt{.tex} file that mirrors the PDF report structure. It:
\begin{itemize}
    \item Generates PNG plot images alongside the \texttt{.tex} file in an \texttt{images/} directory tree.
    \item Global trend plots (\texttt{gain\_vs\_freq.png}, \texttt{bw\_vs\_freq.png}, \texttt{overlay\_cuts.png}) are saved directly in \texttt{images/}.
    \item Per-frequency plots are organised into subfolders: \texttt{images/heatmaps/}, \texttt{images/polar/}, \texttt{images/3d\_surface/}, and \texttt{images/xpol/}.
    \item Uses a well-structured \LaTeX\ document with \texttt{booktabs}, \texttt{longtable}, \texttt{graphicx}, and \texttt{hyperref} packages.
    \item Supports all 13 sections and classification markings.
    \item The generated \texttt{.tex} file can be compiled separately with any standard \LaTeX\ distribution.
\end{itemize}

\subsection{PowerPoint Reports}
The PowerPoint report generator (\texttt{reporting/pptx\_report.py}) produces a \texttt{.pptx} slide deck. It requires the \texttt{python-pptx} package (optional dependency). Key features:
\begin{itemize}
    \item Title slide with project metadata and classification.
    \item Summary table slide with per-frequency statistics.
    \item Trend plot slides (gain vs frequency, beamwidth vs frequency).
    \item Per-frequency heatmap and polar pattern slides.
    \item Configurable via the same 13-section report configuration and 4 presets.
\end{itemize}

\begin{lstlisting}[language=Python,caption={PowerPoint report generation}]
from reporting.pptx_report import generate_pptx_report

generate_pptx_report('./pattern_dir', report_cfg=cfg)
\end{lstlisting}

% ===================================================================
%  PART XII — SIDELOBE MASK STANDARDS
% ===================================================================
\newpage
\section{Sidelobe Mask Standards}
\label{sec:masks}

The sidelobe mask module (\texttt{core/sidelobe\_masks.py}) implements three industry-standard sidelobe envelopes and a custom mask loader.

\subsection{ITU-R S.580-6}
The ITU-R Recommendation S.580-6 defines the reference radiation pattern for earth station antennas used in coordination and interference assessment. The mask is a function of peak gain $G_{\max}$:
\begin{equation}
G(\theta) = \begin{cases}
G_{\max} - 2.5 \times 10^{-3}\left(\frac{D}{\lambda}\theta\right)^2 & \text{for } 0 \leq \theta < \theta_m \\
G_1 & \text{for } \theta_m \leq \theta < \theta_r \\
32 - 25\log_{10}(\theta) & \text{for } \theta_r \leq \theta < 48^\circ \\
-10 & \text{for } 48^\circ \leq \theta \leq 180^\circ
\end{cases}
\end{equation}

\subsection{ITU-R S.465-6}
ITU-R Recommendation S.465-6 is used for coordination of small earth station antennas and is parameterised by $D/\lambda$ (the ratio of antenna diameter to wavelength).

\subsection{MIL-STD}
A military standard sidelobe envelope parameterised by the peak gain and first sidelobe level. Suitable for defence and government applications.

\subsection{Custom Masks}
Custom masks can be loaded from JSON files using \texttt{load\_mask\_from\_json()}. The JSON format defines angle-gain pairs that are linearly interpolated.

\subsection{Compliance Testing}
The \texttt{test\_pattern\_against\_mask()} function:
\begin{enumerate}
    \item Loads the pattern CSV and extracts the selected principal-plane cut.
    \item Evaluates the mask function at each angle.
    \item Compares the pattern gain to the mask limit.
    \item Reports: overall PASS/FAIL, number of violations, violation angles and magnitudes, and minimum margin.
\end{enumerate}

% ===================================================================
%  PART XIII — BATCH PROCESSING
% ===================================================================
\newpage
\section{Batch Processing}
\label{sec:batch_detail}

The batch system (\texttt{core/batch.py}) enables automated generation of multiple antenna configurations from a single JSON recipe file.

\subsection{Recipe Format}
A recipe file is a JSON array of configuration dictionaries. Each entry may contain any subset of \texttt{DEFAULT\_CONFIG} keys — missing keys inherit from defaults. Antenna-specific defaults are automatically merged based on the \texttt{antenna\_type} field.

\subsection{Execution}
\begin{itemize}
    \item \texttt{load\_recipe(filepath)} — Loads and validates the recipe, merging each entry with defaults.
    \item \texttt{run\_batch(recipe\_path, dry\_run, generate\_images, output\_base, writer)} — Iterates over all entries, running \texttt{run\_generation()} for each. Supports dry-run mode for validation without file output. When \texttt{generate\_images=True}, plot images are generated for each job. When \texttt{output\_base} is set, all job output directories are rebased as subfolders under the specified path.
    \item \texttt{generate\_batch\_images(output\_dir, writer)} — Scans a batch output directory for CSV files and generates the full set of plot images (heatmaps, polar, 3D surface, cuts, and multi-frequency trend plots) in a \texttt{graphs/} subfolder tree.
    \item \texttt{create\_sample\_recipe(output\_path)} — Generates a starter template with two example configurations (LPDA and Dish).
\end{itemize}

\subsection{CLI Batch Flags}
\begin{itemize}
    \item \texttt{-{}-batch <recipe.json>} — Run a batch recipe.
    \item \texttt{-{}-batch-dry-run} — Validate the recipe without generating files.
    \item \texttt{-{}-batch-images} — Generate plot images for each batch job.
    \item \texttt{-{}-batch-output <dir>} — Base output directory for all batch jobs.
    \item \texttt{-{}-batch-report} — Generate a PDF report for each batch job (requires \texttt{-{}-batch-images}).
    \item \texttt{-{}-batch-report-preset <name>} — Report section preset: \texttt{default}, \texttt{full}, \texttt{minimal}, or \texttt{quick}.
\end{itemize}

\subsection{Dataset Batch Generator}

The utility script \texttt{scripts/generate\_dataset\_batch.py} generates large-scale randomised dataset recipes for ML training, validation, and benchmarking. It creates batch JSON files with per-type parameter ranges and feature probabilities.

\begin{verbatim}
# Interactive mode
python scripts/generate_dataset_batch.py

# Non-interactive mode
python scripts/generate_dataset_batch.py --count 50000 --types LPDA Panel Horn --seed 42

# All types, random slices, custom output
python scripts/generate_dataset_batch.py --count 10000 --n-slices random --out ./batch_scripts
\end{verbatim}

Features:
\begin{itemize}
    \item Covers all 7 antenna types with physically realistic parameter ranges.
    \item Per-type feature enable probabilities (e.g., 70\% frequency-dependent BW for LPDA, 0\% az BW scaling for Omni).
    \item Log-uniform frequency distribution for even decade coverage.
    \item Automatic Testing/Validation split (default 80/20).
    \item Noise disabled by default (\texttt{noise\_range\_db=0.0}) to produce clean datasets suitable for fitting validation.
    \item Input validation with retry loops: prompts re-ask on invalid input rather than crashing, with sensible defaults (count=100, split=80\%, seed=random).
    \item Reproducible via \texttt{-{}-seed}.
    \item Output: a single \texttt{dataset\_batch.json} in the output directory, consumable by the standard batch processing pipeline.
\end{itemize}

% ===================================================================
%  PART XIV — DEPENDENCY MANAGEMENT
% ===================================================================
\newpage
\section{Dependency Management}
\label{sec:deps}

AntennaForge is designed to run with the Python standard library alone, with optional dependencies for enhanced functionality.

\subsection{Optional Dependencies}

\begin{longtable}{L{3cm} L{10.2cm}}
\toprule
\textbf{Package} & \textbf{Features Enabled} \\
\midrule
\endhead
\texttt{numpy} & Vectorised pattern computation (10--100$\times$ speedup), efficient array operations. \\
\texttt{scipy} & Bicubic interpolation, Bessel functions (Dish model), advanced resampling, RegularGridInterpolator. \\
\texttt{matplotlib} & All plotting, PDF report generation, animation. \\
\texttt{Pillow} & Image loading/display in the GUI plot viewer. \\
\texttt{customtkinter} & The entire GUI framework. Required only for GUI mode. \\
\texttt{python-pptx} & PowerPoint (\texttt{.pptx}) report generation. \\
\bottomrule
\end{longtable}

\subsection{Dependency Checking}
The \texttt{core/deps.py} module provides:
\begin{itemize}
    \item \texttt{check\_all()} — Returns a dictionary of package availability.
    \item \texttt{missing\_packages()} — Lists missing optional packages.
    \item \texttt{install\_command()} — Returns the pip command to install all optional packages.
    \item \texttt{print\_status()} — Prints a human-readable status table.
\end{itemize}

The Home page's Dependency Status card uses these functions to display real-time status.

% ===================================================================
%  PART XV — CONFIGURATION FILE REFERENCE
% ===================================================================
\newpage
\section{Complete Configuration Reference}
\label{sec:config_ref}

The following table documents every key in the \texttt{DEFAULT\_CONFIG} dictionary with its type, default value, valid range, and description.

\subsection{Top-Level Parameters}

\begin{longtable}{L{3.5cm} l l l L{4.7cm}}
\toprule
\textbf{Key} & \textbf{Type} & \textbf{Default} & \textbf{Range} & \textbf{Description} \\
\midrule
\endhead
\texttt{antenna\_type} & str & LPDA & * & Antenna model name \\
\texttt{max\_gain\_dbi} & float & 7.0 & $-30$--$60$ & Peak gain (dBi) \\
\texttt{az\_beamwidth\_deg} & float & 65.0 & $1$--$360$ & Az 3\,dB BW ($^\circ$) \\
\texttt{el\_beamwidth\_deg} & float & 70.0 & $1$--$360$ & El 3\,dB BW ($^\circ$) \\
\texttt{ref\_frequency\_mhz} & float & None & $>0$ & Ref frequency (MHz) \\
\texttt{f\_min\_mhz} & float & 30.0 & $0.001$--$10^6$ & Band lower edge (MHz) \\
\texttt{f\_max\_mhz} & float & 300.0 & $0.001$--$10^6$ & Band upper edge (MHz) \\
\texttt{n\_slices} & int & 10 & $1$--$10000$ & Frequency points \\
\texttt{freq\_spacing} & str & linear & linear, log & Freq distribution \\
\texttt{az\_step\_deg} & float & 1.0 & $0.01$--$45$ & Az angular step ($^\circ$) \\
\texttt{el\_step\_deg} & float & 1.0 & $0.01$--$45$ & El angular step ($^\circ$) \\
\texttt{ftb\_ratio\_db} & float & 15.0 & $0$--$80$ & Front-to-back (dB) \\
\texttt{polarization} & str & vertical & v, h, custom & Polarisation mode \\
\texttt{pol\_angle\_deg} & float & 0.0 & --- & Custom pol angle ($^\circ$) \\
\texttt{noise\_range\_db} & float & 0.0 & $0$--$30$ & Noise amplitude (dB) \\
\texttt{noise\_type} & str & uniform & * & Noise mode: uniform, gaussian, quantization \\
\texttt{sigmoid\_k} & float & 40.0 & $1$--$100$ & Front-fade steepness \\
\texttt{output\_dir} & str & ./antenna\_patterns & --- & Output directory \\
\bottomrule
\end{longtable}

\subsection{Feature Sub-Dictionaries}

Each feature follows the pattern:
\begin{lstlisting}[language={},caption={Feature sub-dictionary structure}]
"feature_name": {
    "enabled": true/false,
    "description": "Human-readable description",
    ...feature-specific parameters...
}
\end{lstlisting}

Refer to Section~\ref{sec:features} for complete documentation of each feature's parameters.

% ===================================================================
%  PART XVI — TIPS AND BEST PRACTICES
% ===================================================================
\newpage
\section{Tips and Best Practices}
\label{sec:tips}

\subsection{Slider and Entry Synchronization}
The \texttt{LabeledSliderEntry} widget maintains bidirectional synchronization between its slider and text entry. Moving the slider updates the entry in real-time; typing in the entry and pressing Enter or clicking elsewhere updates the slider. This allows both coarse adjustment via the slider and precise numeric entry.

\subsection{Tooltips}
Many fields have tooltips that appear on hover after a 400\,ms delay. Tooltips provide concise descriptions and valid ranges for each parameter.

\subsection{Plot Caching}
The Visualize page caches generated plot images. When navigating between plots using the Prev/Next buttons, cached images are displayed instantly without regeneration. The cache is cleared when new plots are generated.

\subsection{Cooperative Cancellation}
Long-running operations (batch plotting, report generation) check a cancellation flag between iterations. Clicking the Cancel button sets this flag, and the operation will stop at the next safe checkpoint.

\subsection{Performance Optimization}
\begin{itemize}
    \item Install NumPy and SciPy for 10--100$\times$ faster pattern generation.
    \item Use coarser angular steps ($2^\circ$ or $5^\circ$) for quick previews; use $1^\circ$ or $0.5^\circ$ for final output.
    \item Reduce the number of frequency slices for initial exploration.
    \item Disable unnecessary physics features (ground reflection, cross-pol) during early design iterations.
\end{itemize}

\subsection{Frequency Spacing}
\begin{itemize}
    \item \textbf{Linear spacing} distributes frequency points evenly across the band. Best for narrowband antennas or when uniform frequency coverage is needed.
    \item \textbf{Logarithmic spacing} places more points at lower frequencies. Ideal for wideband antennas (e.g., LPDAs) where relative bandwidth is more meaningful than absolute bandwidth.
\end{itemize}

\subsection{Reference Frequency}
If left blank, the reference frequency defaults to the geometric mean of the band: $f_{\text{ref}} = \sqrt{f_{\min} \cdot f_{\max}}$. The reference frequency is the frequency at which the specified beamwidth values apply. At other frequencies, the beamwidth is scaled according to the selected decay mode.

\subsection{Output Organisation}
Every generation run\textemdash{}including the very first\textemdash{}creates a timestamped subdirectory (format \texttt{yyyymmddhhmm}) within the output directory. This ensures that no run ever overwrites previous results. The subdirectory name includes the antenna type and the timestamp.

\subsection{Configuration Portability}
Save your generation settings to a JSON file via the Settings page for reproducibility. These configuration files can be:
\begin{itemize}
    \item Shared with colleagues.
    \item Version-controlled alongside your project.
    \item Used as entries in batch recipe files.
    \item Loaded back into the GUI or CLI at any time.
\end{itemize}

% ===================================================================
%  APPENDIX
% ===================================================================
\appendix

\newpage
\section{Default Configuration (JSON)}
\label{app:config}

\begin{lstlisting}[language={},caption={DEFAULT\_CONFIG top-level keys (features omitted for brevity)}]
{
    "antenna_type": "LPDA",
    "max_gain_dbi": 7.0,
    "az_beamwidth_deg": 65.0,
    "el_beamwidth_deg": 70.0,
    "ref_frequency_mhz": null,
    "f_min_mhz": 30.0,
    "f_max_mhz": 300.0,
    "n_slices": 10,
    "freq_spacing": "linear",
    "az_step_deg": 1.0,
    "el_step_deg": 1.0,
    "ftb_ratio_db": 15.0,
    "polarization": "vertical",
    "pol_angle_deg": 0.0,
    "noise_range_db": 0.0,
    "noise_type": "uniform",
    "sigmoid_k": 40.0,
    "output_dir": "./antenna_patterns",
    "features": { ... }
}
\end{lstlisting}

Each feature sub-dictionary follows the structure documented in Section~\ref{sec:features}. Use the ``Show Defaults'' button on the Settings page to view the full configuration at any time.

\newpage
\section{Glossary}
\label{app:glossary}

\begin{longtable}{L{4cm} L{9.2cm}}
\toprule
\textbf{Term} & \textbf{Definition} \\
\midrule
\endhead
\textbf{Azimuth (Az)} & Horizontal angle measured from boresight (front of antenna), $0^\circ$--$360^\circ$. \\
\textbf{Elevation (El)} & Vertical angle measured from the horizon, $-90^\circ$ (nadir) to $+90^\circ$ (zenith). \\
\textbf{Boresight} & The direction of maximum gain, typically $(\mathrm{az}=0^\circ, \mathrm{el}=0^\circ)$. \\
\textbf{Beamwidth (BW)} & The angular width between the $-3$\,dB points of the main beam. \\
\textbf{dBi} & Decibels relative to an isotropic radiator. \\
\textbf{dBm} & Decibels relative to 1 milliwatt. \\
\textbf{EIRP} & Effective Isotropic Radiated Power — the power that would need to be radiated by an isotropic antenna to produce the same power density. \\
\textbf{FSPL} & Free-Space Path Loss — the attenuation of radio waves as they propagate through free space. \\
\textbf{FTB} & Front-to-Back ratio — the gain difference between the front ($0^\circ$) and back ($180^\circ$) of the antenna. \\
\textbf{Sidelobe} & A radiation lobe other than the main beam. \\
\textbf{SLL} & Sidelobe Level — the gain of a sidelobe relative to the main beam peak. \\
\textbf{VSWR} & Voltage Standing Wave Ratio — a measure of impedance mismatch. \\
\textbf{Taylor taper} & An amplitude weighting applied to array elements or used as a sidelobe model, designed to achieve a specified sidelobe level. \\
\textbf{Cross-pol} & Cross-polarisation — the component of radiation in the orthogonal polarisation plane. \\
\textbf{Jinc function} & $\operatorname{jinc}(u) = 2J_1(u)/u$, the Airy pattern for circular apertures. \\
\textbf{Array Factor} & The radiation pattern contribution from the spatial arrangement of array elements. \\
\textbf{Pattern Multiplication} & Total array pattern = element pattern $\times$ array factor. \\
\textbf{Grating Lobe} & Unwanted secondary main beam that appears when element spacing exceeds $\lambda/2$. \\
\bottomrule
\end{longtable}

% ===================================================================
\end{document}
